% Using social networks and collaborative environments — ICT Upper Sixth — Cours signé OnBuch+
% Official MINESEC syllabus. Compiler avec : tectonic ICT12-using-social-networks-and-collaborative-environments.tex
\newcommand{\DOCMATIERE}{ICT}
\newcommand{\DOCNIVEAU}{Upper Sixth}
\newcommand{\DOCMODULE}{Upper Sixth — ICT}
\newcommand{\DOCLECON}{12}
\newcommand{\DOCTITRE}{Using social networks and collaborative environments}
% OnBuch+ — préambule « cours signés » ENRICHI (compilé avec tectonic / XeLaTeX)
% Système visuel « Pop » d'OnBuch+ : fond crème (jamais blanc), bordures encre
% épaisses, ombres « dures » décalées, titres Archivo Black en capitales.
% Version MATHÉMATIQUES : graphes (pgfplots/tikz), boîtes pédagogiques
% (définition, propriété, méthode, attention, le savais-tu, exercice résolu,
% exercice, TP, bilan), page de garde et signature OnBuch+ ancrée en bas.
%
% Macros attendues dans main.tex AVANT \input{preamble} :
%   \DOCMATIERE  \DOCNIVEAU  \DOCMODULE  \DOCLECON (numéro)  \DOCTITRE
\documentclass[11pt]{article}

\usepackage{fontspec}
\usepackage[english]{babel}
\usepackage[a4paper,margin=2cm,top=3.4cm,bottom=2.8cm,headheight=1.4cm,footskip=1.3cm]{geometry}
\usepackage{amsmath,amssymb,amsfonts}
\usepackage{mathtools} % \xmapsto, \coloneqq... souvent utilisés par les modèles
\usepackage{cancel} % simplifications barrées (bilans de forces)
% \abs{z} (module, valeur absolue) et \norm{u} (norme), utilisés par les modèles
\providecommand{\abs}{}\let\abs\relax\DeclarePairedDelimiter{\abs}{\lvert}{\rvert}
\providecommand{\norm}{}\let\norm\relax\DeclarePairedDelimiter{\norm}{\lVert}{\rVert}
\usepackage[table]{xcolor} % [table] : fournit aussi \rowcolors (alternance de lignes)
\usepackage{pagecolor}
\usepackage{tikz}
\usetikzlibrary{shapes.symbols,circuits.logic.US,circuits.logic.IEC,arrows.meta,positioning,calc,shapes.geometric,shapes.misc,shapes.arrows,decorations.pathmorphing,decorations.pathreplacing,decorations.markings,patterns,shadows,mindmap,trees,fit,backgrounds,matrix,intersections,angles,quotes}
\usepackage{pgfplots}
\usepackage[version=4]{mhchem} % équations nucléaires \ce{^{235}_{92}U}
\usepackage[european,siunitx]{circuitikz} % schémas électriques
\pgfplotsset{compat=1.18}
\usepgfplotslibrary{fillbetween,groupplots,polar,statistics,dateplot} % aires entre courbes (calcul intégral)
\usepackage{enumitem}
\usepackage{fancyhdr}
% [most] charge toutes les bibliothèques tcolorbox (listings, minted, raster,
% documentation...) et gonfle inutilement la mémoire TeX sur un document long
% (~300 boîtes) : on ne charge que ce qui est réellement utilisé.
\usepackage[skins,breakable]{tcolorbox}
\usepackage{graphicx}
\usepackage{colortbl}
\usepackage{booktabs}
\usepackage{tabularx}
\usepackage{diagbox} % case d'en-tête diagonale des tableaux à double entrée
% Colonnes extensibles centrée / gauche / droite (C, L, R), souvent utilisées par les modèles
\newcolumntype{C}{>{\centering\arraybackslash}X}
\newcolumntype{L}{>{\raggedright\arraybackslash}X}
\newcolumntype{R}{>{\raggedleft\arraybackslash}X}
\usepackage{array}
\usepackage{multicol}
\usepackage{siunitx}
% parse-numbers=false : accepte aussi \SI{2,5 \times 10^{-3}}{...}
\sisetup{locale=UK,output-decimal-marker={.},group-minimum-digits=4,parse-numbers=false}
\DeclareSIUnit\ohm{\ensuremath{\symup{\Omega}}} % U+2126 absent de la police
\DeclareSIUnit\division{div} % réglages d'oscilloscope (ms/div, V/div)
\DeclareSIUnit\year{yr}\DeclareSIUnit\annum{yr}\DeclareSIUnit\Ma{Ma}\DeclareSIUnit\tr{tr} % tours (vitesse de rotation, ex. \SI{1500}{\tr\per\minute})
\usepackage{qrcode}
% \degree : commande fréquemment utilisée par les modèles pour un angle ou une
% température en dehors de \SI{}{} (non fournie par les paquets chargés).
\providecommand{\degree}{^{\circ}}
% \qed (fin de démonstration), utilisé par les modèles sans amsthm
\providecommand{\qed}{\hfill$\square$}
% \barre : double barre des valeurs interdites dans un tableau de variations (tkz-tab)
\providecommand{\barre}{\ensuremath{\|}}
% environnement proof (amsthm non chargé) utilisé par les modèles
\ifdefined\proof\else\newenvironment{proof}[1][Proof]{\par\noindent\textit{#1.}\ }{\qed\par}\fi
\usepackage{lastpage}
\usepackage{hyperref}
\hypersetup{hidelinks}

% ============================================================
% Palette « Pop » OnBuch+ — mêmes hex que la classe P (plus_ui.dart)
% ============================================================
\definecolor{poporange}{HTML}{F59321}
\definecolor{popdark}{HTML}{E07A0C}
\definecolor{popcream}{HTML}{FFF6E8}
\definecolor{popink}{HTML}{1C1714}
\definecolor{poppurple}{HTML}{7A5AE0}
\definecolor{popgreen}{HTML}{2E9E6A}
\definecolor{poppink}{HTML}{E0457B}
\definecolor{popblue}{HTML}{2F7DE1}
\definecolor{popgold}{HTML}{C98A12}
\definecolor{popmuted}{HTML}{8A7F76}
\definecolor{popline}{HTML}{EADFD2}
\definecolor{popgray}{HTML}{8A7F76}
\colorlet{popgrey}{popgray}
% Alias tolérants (le modèle nomme parfois les couleurs différemment)
\colorlet{popwhite}{white}
\colorlet{popblack}{popink}
\colorlet{popred}{poppink}
\colorlet{popyellow}{popgold}
% Teintes claires pour fonds de tableaux / zones
\colorlet{popblueL}{popblue!10!white}
\colorlet{poporangeL}{poporange!14!white}
\colorlet{popgreenL}{popgreen!12!white}
\colorlet{poppurpleL}{poppurple!12!white}
\colorlet{poppinkL}{poppink!12!white}
\colorlet{popgoldL}{popgold!14!white}

\pagecolor{popcream}

% ============================================================
% Typographie
% ============================================================
\defaultfontfeatures{Ligatures=TeX}
\setmainfont{PlusJakartaSans-Medium.ttf}[Path=../../../fonts/,BoldFont=PlusJakartaSans-Bold.ttf]
% Maths en sans-serif (Fira Math) avec chiffres et lettres droites de Plus
% Jakarta, pour que les formules se fondent dans le texte.
\usepackage[math-style=ISO,bold-style=ISO]{unicode-math}
\setmathfont{FiraMath-Regular.otf}[Path=../../../fonts/]
\setmathfont{PlusJakartaSans-Medium.ttf}[Path=../../../fonts/,range={up/num,up/{Latin,latin},\mathup}]
% (icomma retiré : décimales avec point en anglais)
% Caractères mathématiques Unicode absents des polices de texte (Plus Jakarta,
% Archivo Black) : rendus par leur équivalent mathématique, en texte comme en
% formule — sinon ils s'affichent en carré vide (ex. « Arithmétique dans ℤ »).
\usepackage{newunicodechar}
\newunicodechar{ℝ}{\ensuremath{\mathbb{R}}}
\newunicodechar{ℤ}{\ensuremath{\mathbb{Z}}}
\newunicodechar{ℂ}{\ensuremath{\mathbb{C}}}
\newunicodechar{ℕ}{\ensuremath{\mathbb{N}}}
\newunicodechar{ℚ}{\ensuremath{\mathbb{Q}}}
\newunicodechar{𝔻}{\ensuremath{\mathbb{D}}}
\newunicodechar{α}{\ensuremath{\alpha}}
\newunicodechar{β}{\ensuremath{\beta}}
\newunicodechar{γ}{\ensuremath{\gamma}}
\newunicodechar{δ}{\ensuremath{\delta}}
\newunicodechar{ε}{\ensuremath{\varepsilon}}
\newunicodechar{θ}{\ensuremath{\theta}}
\newunicodechar{λ}{\ensuremath{\lambda}}
\newunicodechar{μ}{\ensuremath{\mu}}
\newunicodechar{σ}{\ensuremath{\sigma}}
\newunicodechar{τ}{\ensuremath{\tau}}
\newunicodechar{φ}{\ensuremath{\varphi}}
\newunicodechar{ψ}{\ensuremath{\psi}}
\newunicodechar{ω}{\ensuremath{\omega}}
\newunicodechar{Σ}{\ensuremath{\Sigma}}
% majuscules produites par \MakeUppercase dans les titres (α-aminés -> Α)
\newunicodechar{Α}{\ensuremath{\alpha}}
\newunicodechar{Β}{\ensuremath{\beta}}
\newunicodechar{Γ}{\ensuremath{\Gamma}}
\newunicodechar{Δ}{\ensuremath{\Delta}}
\newunicodechar{Ε}{\ensuremath{\varepsilon}}
\newunicodechar{Θ}{\ensuremath{\Theta}}
\newunicodechar{Λ}{\ensuremath{\Lambda}}
\newunicodechar{Μ}{\ensuremath{\mu}}
\newunicodechar{Τ}{\ensuremath{\tau}}
\newunicodechar{Φ}{\ensuremath{\Phi}}
\newunicodechar{Ψ}{\ensuremath{\Psi}}
\newunicodechar{Ω}{\ensuremath{\Omega}}
\newunicodechar{↦}{\ensuremath{\mapsto}}
\newunicodechar{∈}{\ensuremath{\in}}
\newunicodechar{∉}{\ensuremath{\notin}}
\newunicodechar{∧}{\ensuremath{\wedge}}
\newunicodechar{⇔}{\ensuremath{\Leftrightarrow}}
\newunicodechar{⟺}{\ensuremath{\Longleftrightarrow}}
\newunicodechar{⇒}{\ensuremath{\Rightarrow}}
\newunicodechar{⟹}{\ensuremath{\Longrightarrow}}
\newunicodechar{∘}{\ensuremath{\circ}}
\newunicodechar{∪}{\ensuremath{\cup}}
\newunicodechar{∩}{\ensuremath{\cap}}
\newunicodechar{≡}{\ensuremath{\equiv}}
\newunicodechar{⊂}{\ensuremath{\subset}}
\newunicodechar{⊥}{\ensuremath{\perp}}
\newunicodechar{∃}{\ensuremath{\exists}}
\newunicodechar{∀}{\ensuremath{\forall}}
\newunicodechar{∛}{\ensuremath{\sqrt[3]{\,}}}
\newunicodechar{∜}{\ensuremath{\sqrt[4]{\,}}}
\newunicodechar{⃗}{\ensuremath{{}^{\to}}}
\newunicodechar{ⁿ}{\ifmmode^{n}\else\textsuperscript{\ensuremath{n}}\fi}
\newunicodechar{ˣ}{\ifmmode^{x}\else\textsuperscript{\ensuremath{x}}\fi}
\newunicodechar{ᵇ}{\ifmmode^{b}\else\textsuperscript{\ensuremath{b}}\fi}
\newunicodechar{ʳ}{\ifmmode^{r}\else\textsuperscript{\ensuremath{r}}\fi}
\newunicodechar{ᵘ}{\ifmmode^{u}\else\textsuperscript{\ensuremath{u}}\fi}
\newunicodechar{ʸ}{\ifmmode^{y}\else\textsuperscript{\ensuremath{y}}\fi}
\newunicodechar{ᵃ}{\ifmmode^{a}\else\textsuperscript{\ensuremath{a}}\fi}
\newunicodechar{ᵉ}{\ifmmode^{e}\else\textsuperscript{\ensuremath{e}}\fi}
\newunicodechar{ᵏ}{\ifmmode^{k}\else\textsuperscript{\ensuremath{k}}\fi}
\newunicodechar{ᵖ}{\ifmmode^{p}\else\textsuperscript{\ensuremath{p}}\fi}
\newunicodechar{ᴺ}{\ifmmode^{N}\else\textsuperscript{\ensuremath{N}}\fi}
\newunicodechar{ᶜ}{\ifmmode^{c}\else\textsuperscript{\ensuremath{c}}\fi}
\newunicodechar{ʰ}{\ifmmode^{h}\else\textsuperscript{\ensuremath{h}}\fi}
\newunicodechar{ᵠ}{\ifmmode^{q}\else\textsuperscript{\ensuremath{q}}\fi}
\newunicodechar{ₙ}{\ifmmode_{n}\else\textsubscript{\ensuremath{n}}\fi}
\newunicodechar{ₐ}{\ifmmode_{a}\else\textsubscript{\ensuremath{a}}\fi}
\newunicodechar{ᵢ}{\ifmmode_{i}\else\textsubscript{\ensuremath{i}}\fi}
\newunicodechar{ᵣ}{\ifmmode_{r}\else\textsubscript{\ensuremath{r}}\fi}
\newunicodechar{ₖ}{\ifmmode_{k}\else\textsubscript{\ensuremath{k}}\fi}
\newunicodechar{ᵦ}{\ifmmode_{\beta}\else\textsubscript{\ensuremath{\beta}}\fi}
\newunicodechar{ₓ}{\ifmmode_{x}\else\textsubscript{\ensuremath{x}}\fi}
\newfontfamily\popdisplay{ArchivoBlack-Regular.ttf}[Path=../../../fonts/]
\newfontfamily\poplogo{SpaceGrotesk-Bold.ttf}[Path=../../../fonts/]
\newfontfamily\popsemi{PlusJakartaSans-SemiBold.ttf}[Path=../../../fonts/]
\newfontfamily\popxbold{PlusJakartaSans-ExtraBold.ttf}[Path=../../../fonts/]
\newfontfamily\popxboldls{PlusJakartaSans-ExtraBold.ttf}[Path=../../../fonts/,LetterSpace=3.0]

\setlist[enumerate]{leftmargin=1.7em,itemsep=5pt,topsep=4pt}
\setlist[itemize]{leftmargin=1.7em,itemsep=4pt,topsep=4pt,label={\color{poporange}\textbullet}}
\setlength{\parindent}{0pt}
\setlength{\parskip}{5pt}
\renewcommand{\arraystretch}{1.25}

% pgfplots : style Pop par défaut
\pgfplotsset{
  every axis/.append style={axis line style={popink,line width=1pt},tick style={popink},
    grid style={popline,line width=0.5pt},label style={font=\small},tick label style={font=\footnotesize}},
  popcurve/.style={poporange,line width=1.8pt,no markers,smooth},
}
% Même style disponible dans un \draw TikZ simple (hors pgfplots)
\tikzset{popcurve/.style={poporange,line width=1.8pt}}
\tikzset{popgrid/.style={popline,very thin}}
\colorlet{popcurve}{poporange} % le nom du style est parfois utilisé comme couleur

% ============================================================
% Pill / squiggle
% ============================================================
\newcommand{\poppill}[3]{%
  \tikz[baseline=-0.55ex]\node[draw=popink,line width=1.2pt,rounded corners=2.6pt,%
    fill=#2,inner xsep=6.5pt,inner ysep=3pt]{%
    \popxboldls\fontsize{7}{7}\selectfont\color{#3}\MakeUppercase{#1}};%
}
\newcommand{\popsquiggle}[1]{%
  \tikz[baseline=-0.4ex]{
    \draw[#1,line width=2.6pt,line cap=round]
      plot [smooth] coordinates {(0,0.09) (0.34,0.01) (0.68,0.21) (1.02,0.01) (1.36,0.10)};
  }%
}
% Mot-clé surligné (marqueur orange sous le texte)
\newcommand{\cle}[1]{{\popxbold\color{popdark}#1}}

% ============================================================
% En-tête / pied
% ============================================================
\pagestyle{fancy}
\fancyhf{}
\renewcommand{\headrulewidth}{0pt}
\renewcommand{\footrulewidth}{0pt}
\lhead{\raisebox{-0.15ex}{\poplogo\fontsize{13}{13}\selectfont\color{popink}OnBuch\color{poporange}+}}
\rhead{\poppill{\DOCMATIERE\ \textbullet\ \DOCNIVEAU\ \textbullet\ Chapter \DOCLECON}{white}{popink}}
\lfoot{\popxbold\fontsize{7.6}{7.6}\selectfont\color{popmuted}\MakeUppercase{Signed course OnBuch+}\ \textbullet\ {\popsemi\DOCTITRE}}
\rfoot{\tikz[baseline=-0.6ex]\node[rounded corners=8.5pt,fill=popink,minimum height=17pt,minimum width=17pt,inner xsep=5pt,inner sep=0]{\popxbold\fontsize{8}{8}\selectfont\color{popcream}\thepage\,/\,\pageref*{LastPage}};}

% ============================================================
% Titres
% ============================================================
\newcommand{\chaptitre}[2]{%
  \begin{tcolorbox}[enhanced,colback=poporange,colframe=popink,boxrule=2.6pt,arc=11pt,
      drop shadow=popink,left=15pt,right=15pt,top=12pt,bottom=11pt,before skip=3pt,after skip=18pt]
    \poppill{#2}{popink}{popcream}\par\vspace{9pt}
    {\raggedright\hyphenpenalty=10000\exhyphenpenalty=10000\popdisplay\fontsize{22}{24}\selectfont\color{popcream}\MakeUppercase{#1}\par}
  \end{tcolorbox}
}
% Section numérotée : pastille orange avec le numéro + titre Archivo
\newcounter{popsec}
\newcommand{\coursec}[1]{%
  \stepcounter{popsec}\setcounter{popsub}{0}%
  \par\vspace{16pt}\noindent
  \begin{minipage}{\linewidth}
  \tikz[baseline=-0.7ex]\node[draw=popink,line width=1.6pt,fill=poporange,rounded corners=4pt,minimum size=22pt,inner sep=0]{\popdisplay\fontsize{12}{12}\selectfont\color{popcream}\thepopsec};%
  \hspace{8pt}{\popdisplay\fontsize{15}{17}\selectfont\color{popink}\MakeUppercase{#1}}\par
  \vspace{4pt}\noindent{\color{poporange}\rule{36pt}{3.6pt}}
  \end{minipage}\par\nopagebreak\vspace{8pt}
}
\newcounter{popsub}
\newcommand{\courssub}[1]{%
  \stepcounter{popsub}%
  \par\vspace{9pt}\noindent{\popxbold\fontsize{11.5}{13}\selectfont\color{popink}{\color{poporange}\thepopsec.\thepopsub}\hspace{6pt}#1}\par\nopagebreak\vspace{4pt}
}

% ============================================================
% Boîtes pédagogiques — même recette : fond blanc, bordure épaisse colorée,
% ombre dure encre, badge plein. Titre optionnel [#1].
% ============================================================
\tcbset{popbase/.style={breakable,enhanced,colback=white,boxrule=2.3pt,arc=9pt,
  shadow={3pt}{-3pt}{0pt}{popink},left=14pt,right=14pt,top=11pt,bottom=11pt,before skip=11pt,after skip=15pt,
  fontupper=\popsemi\fontsize{10.6}{14.5}\selectfont\color{popink},
  fontlower=\popsemi\fontsize{10.6}{14.5}\selectfont\color{popink}}}
% Coche dessinée (le glyphe ✓ n'existe pas dans les polices du document)
\newcommand{\popcheck}[1]{\tikz[baseline=-0.5ex]\draw[#1,line width=1.6pt,line cap=round,line join=round](-0.13,0.0)--(-0.04,-0.09)--(0.13,0.1);}
\providecommand{\coche}{\popcheck{popgreen}}
\providecommand{\resultat}[1]{\par\smallskip\noindent\fcolorbox{popgreen}{popgreenL}{\parbox{\dimexpr\linewidth-2\fboxsep-2\fboxrule\relax}{\textbf{Result.} #1}}\par\smallskip}
\newcommand{\popboxhead}[3]{\poppill{#1}{#2}{white}\ifx\relax#3\relax\else\hspace{6pt}{\popxbold\fontsize{10.6}{12}\selectfont\color{popink}#3}\fi\par\vspace{7pt}}

\newtcolorbox{aretenir}[1][]{popbase,colframe=popgreen,before upper={\popboxhead{Key points}{popgreen}{#1}}}
\newtcolorbox{exemplebox}[1][]{popbase,colframe=poporange,before upper={\popboxhead{Exemple}{poporange}{#1}}}
\newtcolorbox{definition}[1][]{popbase,colframe=popblue,before upper={\popboxhead{Definition}{popblue}{#1}}}
\newtcolorbox{propriete}[1][]{popbase,colframe=poppurple,before upper={\popboxhead{Property}{poppurple}{#1}}}
\newtcolorbox{methode}[1][]{popbase,colframe=popgold,before upper={\popboxhead{Method}{popgold}{#1}}}
\newtcolorbox{attention}[1][]{popbase,colframe=poppink,before upper={\popboxhead{Warning}{poppink}{#1}}}
\newtcolorbox{experience}[1][]{popbase,colframe=popblue,colback=popblueL,before upper={\popboxhead{Experiment}{popblue}{#1}}}
% « Le savais-tu ? » : carte violette pleine (contexte, culture, Cameroun)
\newtcolorbox{savaistu}[1][]{popbase,colback=poppurple,colframe=popink,
  fontupper=\popsemi\fontsize{10.4}{14.2}\selectfont\color{white},
  before upper={\poppill{Did you know?}{white}{poppurple}\ifx\relax#1\relax\else\hspace{6pt}{\popxbold\color{white}#1}\fi\par\vspace{7pt}}}
% Exercice résolu : énoncé en haut, \tcblower puis la solution sur fond crème
\newcounter{popexo}\newcounter{popexr}
\newtcolorbox{exoresolu}[1][]{popbase,colframe=poporange,colbacklower=popcream,
  segmentation style={popink,line width=1.2pt,dashed},
  before upper={\stepcounter{popexr}\popboxhead{Worked example \thepopexr}{poporange}{#1}},
  before lower={\poppill{Solution}{popink}{popcream}\par\vspace{6pt}}}
% Exercice (non résolu en place) — la correction est dans la partie Corrigés
\newtcolorbox{exercice}[1][]{popbase,colframe=popink,
  before upper={\stepcounter{popexo}\popboxhead{Exercise \thepopexo}{popink}{#1}}}
% Corrigé
\newtcolorbox{corrige}[1][]{popbase,colframe=popgreen,colback=popgreenL,
  before upper={\popboxhead{Solution}{popgreen}{#1}}}
% Objectifs / prérequis / bilan
\newtcolorbox{objectifs}{popbase,colframe=popink,colback=poporangeL,
  before upper={\popboxhead{Chapter objectives}{popink}{}}}
\newtcolorbox{prerequis}{popbase,colframe=popink,colback=popblueL,
  before upper={\popboxhead{Prerequisites}{popblue}{}}}
\newtcolorbox{bilan}{popbase,colframe=popink,colback=white,boxrule=2.8pt,
  before upper={\popboxhead{Fiche bilan}{poporange}{L'essentiel à connaître pour l'examen}}}
% Figure encadrée avec légende
\newtcolorbox{popfigure}[1][]{enhanced,breakable=false,colback=white,colframe=popline,boxrule=1.4pt,arc=8pt,
  left=10pt,right=10pt,top=10pt,bottom=8pt,before skip=10pt,after skip=12pt,halign=center,
  lower separated=false,halign lower=center,
  fontlower=\popsemi\fontsize{9}{11.5}\selectfont\color{popmuted},
  #1}
\newcommand{\legende}[1]{\par\vspace{6pt}{\popsemi\fontsize{9}{11.5}\selectfont\color{popmuted}#1}}

% ============================================================
% Signature OnBuch+
% ============================================================
\newcommand{\poplogoinline}[1]{{\poplogo\fontsize{#1}{#1}\selectfont\color{popink}OnBuch\color{poporange}+}}
\newcommand{\signatureonbuch}{%
  \begin{tcolorbox}[enhanced,colback=popink,colframe=popink,boxrule=2.2pt,arc=10pt,
      drop shadow=poporange,left=14pt,right=14pt,top=10pt,bottom=10pt,before skip=0pt,after skip=0pt]
    \tikz[baseline=-0.7ex]\node[circle,fill=popgreen,draw=popcream,line width=1.3pt,minimum size=22pt,inner sep=0]{\popcheck{white}};%
    \hspace{9pt}\begin{minipage}[c]{0.62\linewidth}
      {\popxbold\fontsize{12}{13}\selectfont\color{popcream}Signed\ \textbullet\ {\poplogo OnBuch\color{poporange}+}}\par\vspace{2pt}
      {\raggedright\popsemi\fontsize{8}{10.5}\selectfont\color{popline}\MakeUppercase{OnBuch+ teaching team}\\\MakeUppercase{Follows the official MINESEC syllabus}\par}
    \end{minipage}\hfill
    {\poplogo\fontsize{22}{22}\selectfont\color{popcream}OnBuch\color{poporange}+}
  \end{tcolorbox}%
}

% ============================================================
% Page de garde
% #1 titre · #2 matière · #3 niveau · #4 module · #5 n° leçon · #6 durée ·
% #7 description / objectifs (texte ou itemize)
% ============================================================
\newcommand{\pagegarde}[7]{%
  \thispagestyle{empty}
  \newgeometry{a4paper,margin=1.7cm,top=1.6cm,bottom=1.4cm}%
  \begin{tikzpicture}[remember picture,overlay]
    \fill[poporange!18] ([xshift=-3.2cm,yshift=-3.6cm]current page.north east) circle (4.6cm);
    \fill[poppurple!14] ([xshift=2.4cm,yshift=7.6cm]current page.south west) circle (3.8cm);
  \end{tikzpicture}%
  {\poplogo\fontsize{30}{30}\selectfont\color{popink}OnBuch\color{poporange}+}\hfill
  \raisebox{0.6ex}{\poppill{Signed course \textbullet\ Premium}{popink}{popcream}}\par
  \vspace{1pt}\popsquiggle{poppurple}\par
  \vspace{8pt}
  \begin{tcolorbox}[enhanced,colback=poporange,colframe=popink,boxrule=2.8pt,arc=13pt,
      drop shadow=popink,left=18pt,right=18pt,top=12pt,bottom=13pt,after skip=10pt]
    \poppill{#2 \textbullet\ #3}{popink}{popcream}\hspace{5pt}\poppill{#4}{white}{popink}\par\vspace{8pt}
    {\popdisplay\fontsize{14}{14}\selectfont\color{popink}CHAPTER #5}\par\vspace{4pt}
    {\raggedright\hyphenpenalty=10000\exhyphenpenalty=10000\popdisplay\fontsize{25}{27}\selectfont\color{popcream}\MakeUppercase{#1}\par}
    \vspace{8pt}
    {\popxbold\fontsize{9.5}{9.5}\selectfont\color{popink}\MakeUppercase{Suggested time: #6}}
  \end{tcolorbox}
  \begin{tcolorbox}[enhanced,colback=white,colframe=popink,boxrule=2.2pt,arc=10pt,
      drop shadow=popink,left=16pt,right=16pt,top=10pt,bottom=10pt,after skip=8pt]
    \poppill{In this chapter}{poppurple}{white}\par\vspace{5pt}
    \popsemi\fontsize{9.3}{12.6}\selectfont\color{popink}#7
  \end{tcolorbox}
  \noindent\begin{minipage}[t]{0.487\linewidth}
    \begin{tcolorbox}[enhanced,colback=popgreen,colframe=popink,boxrule=2.2pt,arc=10pt,
        drop shadow=popink,left=11pt,right=11pt,top=10pt,bottom=10pt,height=3.1cm,valign=center]
      \tcbox[colback=white,colframe=popink,boxrule=1.3pt,arc=5pt,boxsep=0pt,nobeforeafter,
        left=3pt,right=3pt,top=3pt,bottom=3pt,valign=center]{\qrcode[height=1.9cm]{https://play.google.com/store/apps/details?id=com.luvvix.onbuch&hl=en}}%
      \hspace{8pt}\begin{minipage}[c]{0.5\linewidth}
        {\popdisplay\fontsize{11}{12.5}\selectfont\color{white}\MakeUppercase{Download OnBuch}}\par\vspace{4pt}
        {\raggedright\popsemi\fontsize{8.4}{11}\selectfont\color{white}Free \textbullet\ Google Play\\AI tutor, past papers, results\par}
      \end{minipage}
    \end{tcolorbox}
  \end{minipage}\hfill
  \begin{minipage}[t]{0.487\linewidth}
    \begin{tcolorbox}[enhanced,colback=poppurple,colframe=popink,boxrule=2.2pt,arc=10pt,
        drop shadow=popink,left=11pt,right=11pt,top=10pt,bottom=10pt,height=3.1cm,valign=center]
      \tikz[baseline=-0.7ex]\node[circle,fill=white,draw=popink,line width=1.3pt,minimum size=1.6cm,inner sep=0]{\popdisplay\fontsize{24}{24}\selectfont\color{poppurple}+};%
      \hspace{8pt}\begin{minipage}[c]{0.6\linewidth}
        {\popdisplay\fontsize{11}{12.5}\selectfont\color{white}\MakeUppercase{Go OnBuch+}}\par\vspace{4pt}
        {\raggedright\popsemi\fontsize{8.4}{11}\selectfont\color{white}All signed courses, unlimited AI tutor, PDF sheets — OnBuch+ tab in the app\par}
      \end{minipage}
    \end{tcolorbox}
  \end{minipage}
  \vspace{8pt}
  \popcontenu
  \vfill
  \signatureonbuch
  \restoregeometry
  \newpage
}

% Rangée « Ce cours contient » (page de garde)
\newcommand{\poptile}[3]{% #1 couleur #2 gros texte #3 légende
  \begin{tcolorbox}[enhanced,colback=#1,colframe=popink,boxrule=2pt,arc=9pt,shadow={2.5pt}{-2.5pt}{0pt}{popink},
      left=6pt,right=6pt,top=9pt,bottom=9pt,halign=center,valign=center,height=1.75cm,width=\linewidth,nobeforeafter]
    {\popdisplay\fontsize{12.5}{13}\selectfont\color{white}\MakeUppercase{#2}}\par\vspace{3pt}
    {\centering\popsemi\fontsize{7.8}{9.5}\selectfont\color{white}#3\par}
  \end{tcolorbox}}
\newcommand{\popcontenu}{%
  \par\noindent{\popxboldls\fontsize{7.5}{7.5}\selectfont\color{popmuted}THIS SIGNED COURSE CONTAINS}\par\vspace{7pt}
  \noindent\begin{minipage}[t]{0.235\linewidth}\poptile{popblue}{Course}{complete, illustrated, step by step}\end{minipage}\hfill
  \begin{minipage}[t]{0.235\linewidth}\poptile{poporange}{Methods}{and worked examples}\end{minipage}\hfill
  \begin{minipage}[t]{0.235\linewidth}\poptile{poppink}{Exercises}{exam style + solutions}\end{minipage}\hfill
  \begin{minipage}[t]{0.235\linewidth}\poptile{popgreen}{Summary sheet}{the essentials to remember}\end{minipage}\par}

% Sommaire
\newtcolorbox{sommaire}{enhanced,colback=white,colframe=popink,boxrule=2.2pt,arc=10pt,
  drop shadow=popink,left=18pt,right=18pt,top=13pt,bottom=13pt,before skip=6pt,after skip=20pt,
  before upper={\poppill{Contents}{poppurple}{white}\par\vspace{9pt}\popsemi\fontsize{10.6}{14.5}\selectfont\color{popink}}}

% Signature de fin de cours — poussée tout en bas de la dernière page
\newcommand{\findecours}{%
  \par\vfill\nopagebreak
  \begin{center}\popsquiggle{poporange}\end{center}\vspace{4pt}
  \signatureonbuch
}
\AtBeginDocument{\def\llbracket{\mathopen{[\mkern-3.5mu[}}\def\rrbracket{\mathclose{]\mkern-3.5mu]}}} % intervalles d'entiers (glyphe ⟦ absent de la police)
% Glyphes absents de Fira Math : redéfinis à partir de glyphes disponibles
\AtBeginDocument{%
  \def\setminus{\mathbin{\backslash}}\let\smallsetminus\setminus
  \def\vdots{\mathord{\vbox{\baselineskip3.2pt\lineskiplimit0pt\kern4pt\hbox{.}\hbox{.}\hbox{.}}}}%
  \def\ddots{\mathinner{\mkern1mu\raise6.5pt\hbox{.}\mkern2mu\raise3.6pt\hbox{.}\mkern2mu\raise0.7pt\hbox{.}\mkern1mu}}%
  \def\triangle{\mathord{\scalebox{0.9}{$\symup{\Delta}$}}}%
  \def\nabla{\mathord{\raisebox{\depth}{\scalebox{1}[-1]{$\symup{\Delta}$}}}}%
  \def\checkmark{\popcheck{popdark}}%
  \def\star{\mathbin{\ast}}\def\bigstar{\mathord{\ast}}%
  \def\diamond{\mathbin{\scalebox{0.6}{$\Diamond$}}}%
  \def\bigcup{\mathop{\mathchoice{\raisebox{-0.3em}{\scalebox{1.7}{$\cup$}}}{\scalebox{1.25}{$\cup$}}{\cup}{\cup}}\displaylimits}%
  \def\bigcap{\mathop{\mathchoice{\raisebox{-0.3em}{\scalebox{1.7}{$\cap$}}}{\scalebox{1.25}{$\cap$}}{\cap}{\cap}}\displaylimits}%
  \def\lesseqqgtr{\mathrel{\lessgtr}}\def\bigoplus{\mathop{\oplus}\displaylimits}%
}
\providecommand{\ssi}{if and only if} % abréviation fréquente
\AtBeginDocument{\providecommand{\wideparen}{\overparen}} % arc de cercle

% Fonctions usuelles que les modèles utilisent sans que amsmath les fournisse
\providecommand{\arsinh}{\operatorname{arsinh}}\providecommand{\arcosh}{\operatorname{arcosh}}
\providecommand{\artanh}{\operatorname{artanh}}\providecommand{\arsech}{\operatorname{arsech}}
\providecommand{\arcsch}{\operatorname{arcsch}}\providecommand{\arcoth}{\operatorname{arcoth}}
\providecommand{\arcsinh}{\operatorname{arsinh}}\providecommand{\arccosh}{\operatorname{arcosh}}
\providecommand{\arctanh}{\operatorname{artanh}}\providecommand{\sech}{\operatorname{sech}}
\providecommand{\csch}{\operatorname{csch}}\providecommand{\arcsec}{\operatorname{arcsec}}
\providecommand{\arccsc}{\operatorname{arccsc}}\providecommand{\arccot}{\operatorname{arccot}}
\providecommand{\sgn}{\operatorname{sgn}}\providecommand{\lcm}{\operatorname{lcm}}
\providecommand{\Var}{\operatorname{Var}}\providecommand{\Cov}{\operatorname{Cov}}
\providecommand{\permil}{\ensuremath{{}^{0}\!/_{00}}}\providecommand{\textperthousand}{\ensuremath{{}^{0}\!/_{00}}}

%% FIGFIX-BEGIN v5 — correctifs globaux des figures (docs/onbuch-generation/tools/figfix)
\usepackage{adjustbox}\usepackage{etoolbox}\tcbuselibrary{raster}
% toute figure TikZ/pgfplots plus large que la ligne est réduite à la largeur disponible
\BeforeBeginEnvironment{tikzpicture}{\begin{adjustbox}{max width=\linewidth}}
\AfterEndEnvironment{tikzpicture}{\end{adjustbox}}
% légendes pgfplots lisibles par-dessus les courbes
\pgfplotsset{every axis/.append style={legend style={fill=white,fill opacity=0.92,text opacity=1,draw=black!15,inner sep=3pt}}}
% étiquettes d'arêtes (syntaxe edge["..."]) avec fond blanc
\tikzset{every edge quotes/.append style={fill=white,inner sep=1.2pt}}
%% FIGFIX-END
\begin{document}

\pagegarde{Using social networks and collaborative environments}{ICT}{Upper Sixth}{Upper Sixth — ICT}{12}{1 periods}{This chapter explores social networks and collaborative environments: their definitions, types, architectures, uses, advantages and risks. Students learn to exploit these tools for learning, communication and teamwork while mastering netiquette, privacy protection and responsible online behaviour expected at GCE Advanced Level.
\vspace{4pt}
\begin{multicols}{2}\small
\begin{itemize}
\item By the end of this chapter I can define a social network and a collaborative environment and distinguish between them
\item By the end of this chapter I can classify social networks (general, professional, media-sharing, messaging, forums) with Cameroonian and international examples
\item By the end of this chapter I can describe the main features of a social network account: profile, friends/followers, news feed, groups, pages, privacy settings
\item By the end of this chapter I can create, configure and secure an account on a social network or collaborative platform
\item By the end of this chapter I can use collaborative tools (Google Workspace, Microsoft 365, wikis, cloud storage, videoconferencing) to produce shared work in a team
\item By the end of this chapter I can explain the advantages of social networks and collaborative work for education, business and civic life in Cameroon
\end{itemize}\end{multicols}}

\begin{sommaire}
\begin{multicols}{2}
\begin{enumerate}
\item Social Networks: Concepts, Types and Key Actors
\item Using a Social Network: Accounts, Features and Privacy
\item Collaborative Environments and Tools for Teamwork
\item Risks, Ethics, Netiquette and Legal Framework
\item Integration activity
\item Methods and worked examples
\item Exercises
\item Solutions to the exercises
\item Summary sheet
\end{enumerate}
\end{multicols}
\end{sommaire}

\begin{objectifs}
\begin{itemize}
\item By the end of this chapter I can define a social network and a collaborative environment and distinguish between them
\item By the end of this chapter I can classify social networks (general, professional, media-sharing, messaging, forums) with Cameroonian and international examples
\item By the end of this chapter I can describe the main features of a social network account: profile, friends/followers, news feed, groups, pages, privacy settings
\item By the end of this chapter I can create, configure and secure an account on a social network or collaborative platform
\item By the end of this chapter I can use collaborative tools (Google Workspace, Microsoft 365, wikis, cloud storage, videoconferencing) to produce shared work in a team
\item By the end of this chapter I can explain the advantages of social networks and collaborative work for education, business and civic life in Cameroon
\item By the end of this chapter I can identify the dangers of social networks (cyberbullying, fake news, scams, addiction, privacy violations) and apply protective measures
\item By the end of this chapter I can apply the rules of netiquette and responsible digital citizenship, including respect of Cameroonian law on cybercrime (Law No. 2010/012)
\end{itemize}
\end{objectifs}

\begin{prerequis}
\begin{itemize}
\item Internet, the Web, browsers and search engines (Lower Sixth)
\item Basic computer networks: clients, servers, IP addresses (Lower Sixth)
\item E-mail and electronic messaging (Form 5 / Lower Sixth)
\item Computer security basics: passwords, malware, antivirus (Lower Sixth)
\item Word processing and file management
\end{itemize}
\end{prerequis}

\newpage

\coursec{Social Networks: Concepts, Types and Key Actors}

Every morning in Yaoundé, students check WhatsApp groups before even leaving bed, traders in Douala's Marché Mboppi advertise their goods on Facebook Marketplace, and a class in Bamenda shares revision notes through a Google Doc edited by twenty students at once. Yet the same week, a candidate's GCE results were ``leaked'' through a fake Facebook page, and a young entrepreneur lost her savings to a scammer met on a social network. Social networks and collaborative environments have become the invisible infrastructure of Cameroonian life --- for learning, business, politics and entertainment. But do we really understand how they work, what they do with our data, and how to use them professionally and safely? This chapter turns everyday users into informed, critical and productive digital citizens.

In this first section, we lay the foundations: what exactly a social network is, where these platforms came from, how they are classified, how they work technically, and why they are ``free''.

\courssub{What is a Social Network?}

\textbf{Intuition first.} Imagine a large hall in your school where every student has a personal notice board on the wall. On your board you pin your photo, your name and your interests. You can walk to a friend's board and pin a message on it, and your friend can pin one on yours. Anyone passing by can see who is connected to whom, and what has been shared. A social network is exactly this hall --- except the hall is a set of powerful computers (servers) somewhere in the world, the boards are \emph{profiles}, and the walking is done by your phone over the Internet.

\begin{definition}[Social network]
A \cle{social network} is an online platform that enables users to (i) create a \cle{public or semi-public profile}, (ii) \cle{connect} with other users (friends, followers, contacts), and (iii) \cle{create and share content} (text, photos, videos, links) with their network, and to interact with the content of others (like, comment, share).
\end{definition}

The three elements of the definition are all essential, and they give us a test: a platform is a social network only if it offers \textbf{profiles}, \textbf{connections between users} and \textbf{content sharing and interaction}. We shall use this test later to distinguish social networks from simple websites, forums and messaging applications.

\begin{attention}[Common mistake: social network $\neq$ search engine $\neq$ Internet]
Candidates frequently confuse three different things in the GCE examination:
\begin{itemize}
\item \textbf{Google} is a \cle{search engine}: it indexes web pages so that you can find them. It does not exist so that users connect to each other.
\item \textbf{Facebook, TikTok, WhatsApp} are \cle{social platforms}: services built \emph{on top of} the Internet.
\item The \textbf{Internet} is the underlying \cle{network infrastructure} (cables, routers, protocols such as TCP/IP) that carries all these services.
\end{itemize}
Saying ``I found it on the Internet'' when you mean ``on Facebook'', or ``I searched on Facebook'' when you mean ``on Google'', costs marks. The Internet is the road; Google is a signpost; Facebook is a marketplace built beside the road.
\end{attention}

\courssub{A Short History of Social Networks}

Social networking did not begin with Facebook. The key milestones every Upper Sixth student should know are:

\begin{center}
\begin{tabularx}{\linewidth}{|l|l|X|}
\hline
\rowcolor{popblueL} \textbf{Year} & \textbf{Platform} & \textbf{Contribution} \\
\hline
1997 & SixDegrees & Often considered the first recognisable social network: profiles, friend lists and browsing of connections. \\
\hline
2004 & Facebook & Launched for university students, then opened to everyone; popularised the real-name profile and the news feed. \\
\hline
2006 & Twitter (now X) & Popularised \cle{microblogging}: very short public messages. \\
\hline
2009 & WhatsApp & Instant messaging built around the phone number; later added groups, status and voice/video calls. \\
\hline
2010 & Instagram & Photo and video sharing centred on images. \\
\hline
2016 & TikTok & Short viral videos recommended by a powerful algorithm. \\
\hline
\end{tabularx}
\end{center}

\begin{savaistu}[Did you know?]
The name ``SixDegrees'' came from the idea of \emph{six degrees of separation}: the claim that any two people on Earth are connected through a short chain of acquaintances. Social networks made this idea visible and measurable, because they store the chain of connections explicitly in their databases.
\end{savaistu}

\courssub{Classification of Social Networks}

Not all social platforms serve the same purpose. To classify a platform, examiners expect you to identify its \emph{dominant purpose} and the \emph{type of content} exchanged.

\begin{methode}[Classifying an unknown platform]
When faced with a platform you do not know, ask two questions:
\begin{enumerate}
\item \textbf{Who connects to whom?} Friends and family? Professionals and employers? Anonymous fans of a topic?
\item \textbf{To share what?} Short text messages? Photos and videos? CVs and job offers? Private chats?
\end{enumerate}
The pair of answers determines the category. For example: ``professionals connecting to employers to share CVs'' $\Rightarrow$ professional network (LinkedIn).
\end{methode}

\begin{center}
\begin{tabularx}{\linewidth}{|l|X|X|}
\hline
\rowcolor{popblueL} \textbf{Category} & \textbf{Purpose / content} & \textbf{Examples} \\
\hline
General networks & Connect friends, family and acquaintances; mixed content & Facebook \\
\hline
Professional networks & Careers, CVs, recruitment, business contacts & LinkedIn \\
\hline
Microblogging & Very short public posts, news and opinions in real time & X (Twitter) \\
\hline
Media sharing & Photos and videos as the main content & YouTube, Instagram, TikTok \\
\hline
Instant messaging & Private or group conversations in real time & WhatsApp, Telegram \\
\hline
Forums and communities & Thematic discussions organised in threads & Reddit, student discussion boards \\
\hline
\end{tabularx}
\end{center}

\begin{popfigure}
\begin{tikzpicture}[
  grow=right,
  level 1/.style={sibling distance=52mm, level distance=34mm},
  level 2/.style={sibling distance=11mm, level distance=30mm},
  every node/.style={font=\small},
  edge from parent/.style={draw=popmuted, thick}
]
\node[draw=popdark, fill=popblueL, rounded corners, thick, align=center, font=\small\bfseries] {Social\\networks}
  child { node[draw=popline, rounded corners, align=center] {Forums /\\communities}
    child { node[draw=popmuted, rounded corners, align=center, font=\scriptsize] {Reddit} }
  }
  child { node[draw=popline, rounded corners, align=center] {Instant\\messaging}
    child { node[draw=popmuted, rounded corners, align=center, font=\scriptsize] {Telegram} }
    child { node[draw=popmuted, rounded corners, align=center, font=\scriptsize] {WhatsApp} }
  }
  child { node[draw=popline, rounded corners, align=center] {Media\\sharing}
    child { node[draw=popmuted, rounded corners, align=center, font=\scriptsize] {TikTok} }
    child { node[draw=popmuted, rounded corners, align=center, font=\scriptsize] {YouTube} }
  }
  child { node[draw=popline, rounded corners, align=center] {Microblogging}
    child { node[draw=popmuted, rounded corners, align=center, font=\scriptsize] {X (Twitter)} }
  }
  child { node[draw=popline, rounded corners, align=center] {Professional}
    child { node[draw=popmuted, rounded corners, align=center, font=\scriptsize] {LinkedIn} }
  }
  child { node[draw=popline, rounded corners, align=center] {General}
    child { node[draw=popmuted, rounded corners, align=center, font=\scriptsize] {Facebook} }
  };
\end{tikzpicture}
\legende{Classification tree of social networks with one or two examples per category. Some platforms (e.g.\ WhatsApp with its ``status'' feature) sit on the border between categories.}
\end{popfigure}

\begin{attention}[Borderline cases]
Real platforms evolve and borrow features from each other: WhatsApp added ``status'' (a media-sharing feature), Facebook added Marketplace (e-commerce). In the exam, always classify a platform by its \textbf{main purpose}, and justify your choice with the two questions of the method above.
\end{attention}

\courssub{Social Networks in Cameroon}

Social networks are deeply woven into Cameroonian life. The most used platforms are typically \textbf{Facebook}, \textbf{WhatsApp}, \textbf{TikTok}, \textbf{YouTube} and \textbf{Instagram}, accessed mostly through smartphones on mobile data bundles offered by operators such as MTN and Orange.

Their uses go far beyond entertainment:

\begin{itemize}
\item \textbf{Business.} Traders advertise goods on \cle{Facebook Marketplace}; small businesses manage customers and catalogues through \cle{WhatsApp Business}; mobile money (MTN MoMo, Orange Money) is often used to close sales agreed on these platforms.
\item \textbf{Education.} Classes create WhatsApp groups to share notes and announcements; teachers post revision videos on YouTube; during school disruptions, many courses continued through these channels.
\item \textbf{Politics and social mobilisation.} Political parties, civil society and citizens debate public issues on Facebook and X; hashtags have been used to mobilise opinion around national events.
\item \textbf{Information --- and misinformation.} News spreads extremely fast, but so do rumours and fake pages (as in the fake GCE results page of our opening story). Critical verification is an essential skill, developed later in this chapter.
\end{itemize}

\begin{popfigure}
\begin{tikzpicture}
\begin{axis}[
  ybar,
  width=0.95\linewidth,
  height=6.2cm,
  bar width=0.55cm,
  ymin=0, ymax=3.4,
  ylabel={Users (billions, approx.)},
  symbolic x coords={Facebook, YouTube, WhatsApp, Instagram, TikTok},
  xtick=data,
  x tick label style={font=\small},
  ytick={0,1,2,3},
  nodes near coords,
  every node near coord/.append style={font=\scriptsize},
  axis lines=left,
  ymajorgrids=true,
  grid style={popline, dashed}
]
\addplot[fill=popblue, draw=popdark] coordinates {
  (Facebook, 3.0)
  (YouTube, 2.5)
  (WhatsApp, 2.4)
  (Instagram, 2.0)
  (TikTok, 1.6)
};
\end{axis}
\end{tikzpicture}
\legende{Approximate monthly active users of major platforms worldwide (orders of magnitude, in billions). Values are indicative and change every year --- what matters for the exam is the \emph{relative} scale, not the exact figures.}
\end{popfigure}

\courssub{How a Social Network Works Technically}

\textbf{Observation.} When you open Facebook on your phone in Buea, the photos of a friend in Bamenda appear in seconds --- even though those photos are not stored on your phone. Where were they? On the platform's \cle{servers}.

\begin{definition}[Client--server model]
In the \cle{client--server model}, a \cle{client} (your web browser or mobile application) sends \cle{requests} over the Internet to a \cle{server} (a powerful computer run by the platform), which processes the request and sends back a \cle{response} (a page, a photo, a list of posts). User accounts, profiles, photos and messages are all stored in \cle{databases} on the platform's servers, not on your device.
\end{definition}

Concretely, when you open the app: (1) your phone authenticates you (login, password); (2) it requests your \cle{news feed}; (3) the server selects posts from its database; (4) it sends them back; (5) the app displays them.

\begin{definition}[News feed and algorithm]
The \cle{news feed} (or timeline) is the continuously updated list of posts displayed to a user. It is \textbf{not} simply chronological: it is generated by an \cle{algorithm}, that is, a finite, ordered sequence of instructions that solves a problem --- here, the problem ``which posts, among millions, should this user see first?''. Feed algorithms rank posts using signals such as your past likes, the people you interact with most, the type of content you watch longest, and the recency of posts.
\end{definition}

Why does this matter? Because the algorithm decides what you see --- and therefore what you \emph{think} is important. Two users in the same WhatsApp group may open Facebook and see completely different ``realities''. Understanding this is the first step towards critical use.

\begin{popfigure}
\begin{tikzpicture}[font=\small, >=stealth]
% Server
\node[draw=popdark, fill=popblueL, rounded corners, thick, minimum width=3.4cm, minimum height=2.2cm, align=center] (srv) at (6,0)
  {\textbf{Platform servers}\\[1pt] \scriptsize accounts, profiles,\\ \scriptsize photos, databases};
% Clients
\node[draw=popdark, fill=popgreenL, rounded corners, thick, minimum width=1.7cm, minimum height=1.3cm, align=center] (phone) at (0,1.8)
  {\textbf{Smartphone}\\ \scriptsize mobile app};
\node[draw=popdark, fill=popgreenL, rounded corners, thick, minimum width=1.7cm, minimum height=1.3cm, align=center] (pc) at (0,-1.8)
  {\textbf{PC}\\ \scriptsize web browser};
% Arrows phone
\draw[->, thick, popblue] (phone.east) to[bend left=10] node[midway, above, font=\scriptsize, align=center] {1.\ request\\ \scriptsize (login, feed)} (srv.160);
\draw[->, thick, poporange] (srv.200) to[bend left=10] node[midway, below, font=\scriptsize, align=center] {2.\ response\\ \scriptsize (news feed)} (phone.east);
% Arrows PC
\draw[->, thick, popblue] (pc.east) to[bend right=10] node[midway, above, font=\scriptsize, align=center] {1.\ request} (srv.340);
\draw[->, thick, poporange] (srv.300) to[bend right=10] node[midway, below, font=\scriptsize, align=center] {2.\ response} (pc.east);
% Internet cloud label
\node[font=\scriptsize\itshape, text=popmuted, align=center] at (3.1,0) {Internet\\ (TCP/IP)};
\end{tikzpicture}
\legende{Client--server operation of a social network: clients (smartphone app, browser) send requests; the servers consult their databases and return the personalised news feed.}
\end{popfigure}

\courssub{The Business Model: ``If It Is Free, You Are the Product''}

\textbf{Observation.} Facebook, TikTok and YouTube cost nothing to join, yet their parent companies are among the richest in the world. Who pays? Not you --- at least not with money.

\begin{aretenir}[Key fact: the advertising business model]
Most social networks are \cle{free services financed by targeted advertising}. The platform collects data about users (age, location, interests, pages liked, videos watched, time spent) and uses it to show each user \emph{personalised advertisements}, which advertisers pay for. Hence the famous saying: \emph{``If it is free, you are the product.''} Your attention and your data are what is being sold. This also explains why feed algorithms favour content that keeps you online longer: more time online means more advertisements seen.
\end{aretenir}

This model has direct consequences for you as a citizen: your personal data has economic value, which is why laws on data protection exist (in Cameroon, this area involves institutions such as ANTIC, the National Agency for Information and Communication Technologies) and why privacy settings --- studied in the next section --- matter so much.

\courssub{Essential Terminology}

The GCE expects precise vocabulary. Learn this table carefully:

\begin{center}
\begin{tabularx}{\linewidth}{|l|X|}
\hline
\rowcolor{popblueL} \textbf{Term} & \textbf{Meaning} \\
\hline
\cle{Profile} & The personal page describing a user: name, photo, information, publications. \\
\hline
\cle{Friend / Follower} & A user connected to you: friendship is usually mutual (Facebook), following is often one-way (X, Instagram). \\
\hline
\cle{Post} & A piece of content published by a user (text, photo, video, link). \\
\hline
\cle{Like} & A one-click positive reaction to a post. \\
\hline
\cle{Comment} & A written public reply attached to a post. \\
\hline
\cle{Share} & Republishing someone else's post to your own network. \\
\hline
\cle{Hashtag} & A keyword preceded by \# (e.g.\ \#GCE2025) that groups posts on the same topic. \\
\hline
\cle{Tag} & Mentioning a user (with @) in a post or photo, which notifies them. \\
\hline
\cle{Story} & Ephemeral content that disappears after a short time (usually 24 hours). \\
\hline
\cle{News feed} & The algorithmically generated stream of posts shown to a user. \\
\hline
\cle{Group} & A community space where members with a common interest post and discuss. \\
\hline
\cle{Page} & A public profile for an organisation, business or public figure, which users can follow. \\
\hline
\end{tabularx}
\end{center}

\courssub{Exam Extra: Social Network, Website, Forum or Messaging App?}

A frequent GCE question asks you to \textbf{distinguish a social network from related tools}. Apply the three-part test from our definition (profile + connections + sharing/interaction):

\begin{center}
\begin{tabularx}{\linewidth}{|l|X|X|}
\hline
\rowcolor{popblueL} \textbf{Tool} & \textbf{Main purpose} & \textbf{Why it is (not) a social network} \\
\hline
Simple website (e.g.\ a school's site) & Publish information to visitors & No profiles, no connections between users $\Rightarrow$ \textbf{not} a social network. \\
\hline
Discussion forum & Thematic discussions in threads & Users interact around \emph{topics}, not around personal networks; often considered an ancestor or borderline form of social network. \\
\hline
Instant messaging (WhatsApp) & Private real-time conversations & Connections exist, but content is private, not published to a network; classified as a messaging platform with social features. \\
\hline
Social network (Facebook) & Profiles + connections + public sharing & All three criteria satisfied. \\
\hline
\end{tabularx}
\end{center}

\begin{exoresolu}[Classify and justify]
\textbf{Statement.} A GCE candidate writes: ``Google is my favourite social network because I find everything there.'' (a) Identify the error. (b) Classify correctly: Google, Facebook, WhatsApp. (c) Justify each classification in one sentence.
\tcblower
\textbf{Solution.}
\begin{enumerate}
\item[(a)] The candidate confuses a \cle{search engine} with a \cle{social network}: finding information is not the same as connecting users who share content.
\item[(b)--(c)] \textbf{Google}: a search engine --- it indexes web pages and returns links; users do not build profiles or connect to each other. \textbf{Facebook}: a general social network --- users create profiles, add friends and share posts publicly or semi-publicly. \textbf{WhatsApp}: an instant messaging platform with social features --- it connects users through phone numbers for private and group conversations rather than public publishing.
\end{enumerate}
\end{exoresolu}

\begin{exercice}[Check your understanding]
\begin{enumerate}
\item Give the three essential elements that define a social network.
\item Using the two-question method, classify LinkedIn and TikTok, justifying each answer.
\item Explain, in three sentences, the meaning of ``If it is free, you are the product.''
\item Draw and label a client--server diagram showing what happens when a student in Limbe opens her news feed.
\item State one difference between a \emph{group} and a \emph{page} on Facebook.
\end{enumerate}
\end{exercice}

\begin{corrige}[Exercise 1]
\begin{enumerate}
\item Public or semi-public \textbf{profiles}; \textbf{connections} between users; \textbf{sharing of content} with interaction (like, comment, share).
\item \textbf{LinkedIn}: professionals connect to employers and colleagues to share CVs and career content $\Rightarrow$ professional network. \textbf{TikTok}: users connect to creators to share and watch short videos $\Rightarrow$ media-sharing network.
\item The service is free because revenue comes from advertisers; the platform collects user data and attention to display targeted advertisements; therefore the user's data and attention are the real merchandise sold to advertisers.
\item The diagram must show: the smartphone (client) sending a request through the Internet to the platform's servers, and the servers returning the personalised news feed generated from their databases (mirror the figure of this section).
\item A \textbf{group} is a community space where members discuss a common interest; a \textbf{page} is a public profile for an organisation, business or public figure that users follow.
\end{enumerate}
\end{corrige}

\begin{aretenir}[Section summary]
\begin{itemize}
\item A social network = \textbf{profiles} + \textbf{connections} + \textbf{content sharing and interaction}.
\item Milestones: SixDegrees (1997), Facebook (2004), Twitter/X (2006), WhatsApp (2009), TikTok (2016).
\item Categories: general, professional, microblogging, media sharing, instant messaging, forums/communities --- classify by \emph{who connects to whom, to share what}.
\item Technically, platforms use the \textbf{client--server model}; the \textbf{news feed} is produced by a ranking \textbf{algorithm}.
\item Free platforms are financed by \textbf{targeted advertising} based on user data.
\item Never confuse a social network with a search engine, a simple website, a forum or the Internet itself.
\end{itemize}
\end{aretenir}

\coursec{Using a Social Network: Accounts, Features and Privacy}

In the previous section we discovered what social networks are, the main types (general networks, professional networks, media-sharing platforms, messaging services) and the key actors that operate them. Knowing the landscape is not enough, however: to benefit from these platforms while avoiding their dangers, an Upper Sixth student must master the \cle{practical skills} of creating and securing an account, managing a profile, publishing wisely and protecting privacy. This section takes you step by step through the complete life of a social network account, from registration to the protection of your digital identity.

\courssub{Creating and Securing an Account}

\textbf{Registration data.}\par
Joining a social network always begins with a \cle{registration form}. The platform typically asks for a name, a date of birth, an e-mail address or a mobile phone number, and a password. Two principles should guide you here. First, provide only what is \emph{required}: fields marked optional (home town, school, relationship status) can usually be left empty. Second, be truthful about essentials such as your age, because most platforms legally forbid children under a minimum age (commonly 13 years), and lying about your identity can make account recovery impossible later.

After registration, the platform asks you to \cle{verify} your e-mail address or phone number, usually by clicking a link or entering a code sent by SMS. This verification proves that you really own the contact address; it is also your lifeline if you ever forget your password.

\textbf{Choosing a strong password.}\par
Your password is the key to your account. A \cle{strong password} is long (at least 12 characters), mixes upper-case and lower-case letters, digits and symbols, and is \emph{unique} to that account. Attackers use two main techniques: \emph{dictionary attacks}, which try common words and names, and \emph{brute-force attacks}, which try every combination. Length is your best defence: each extra character multiplies the number of combinations an attacker must test.

\begin{exemplebox}[Weak and strong passwords]
The password \texttt{ngong2024} is weak: it contains a name and a year, both easy to guess from a public profile. A passphrase such as \texttt{Mboa!Rain3Falls\$Douala} is far stronger: it is long, mixes character types, and is not found in any dictionary. Never reuse your e-mail password on a social network: if one service is breached, all your accounts fall together.
\end{exemplebox}

\textbf{Two-factor authentication (2FA).}\par
Even a strong password can be stolen. \cle{Two-factor authentication} adds a second proof of identity: after entering your password, you must also enter a one-time code (OTP) sent by SMS or generated by an authenticator application (e.g. Google Authenticator, Microsoft Authenticator). The idea is to combine \emph{something you know} (the password) with \emph{something you have} (your phone). A thief who steals your password in Buea cannot log in without the phone in your pocket in Bamenda.

\begin{popfigure}
\begin{center}
\begin{tikzpicture}[
  node distance=0.55cm,
  box/.style={draw=popblue, thick, rounded corners, fill=popblueL, text width=3.6cm, align=center, minimum height=1.2cm, font=\small},
  arr/.style={->, thick, popdark}]
\node[box, align=center] (reg){\textbf{1. Register}\\ name, e-mail or phone, date of birth};
\node[box, right=of reg, align=center] (ver){\textbf{2. Verify}\\ confirm e-mail link or SMS code};
\node[box, right=of ver, align=center] (pwd){\textbf{3. Strong password}\\ 12+ characters, unique, mixed types};
\node[box, below=0.7cm of pwd, align=center] (tfa){\textbf{4. Activate 2FA}\\ OTP by SMS or authenticator app};
\node[box, left=of tfa, align=center] (priv){\textbf{5. Privacy settings}\\ audience, contact, location, apps};
\node[box, left=of priv, fill=popgreenL, draw=popgreen, align=center] (ok){\textbf{6. Secure account}\\ review settings regularly};
\draw[arr] (reg) -- (ver);
\draw[arr] (ver) -- (pwd);
\draw[arr] (pwd) -- (tfa);
\draw[arr] (tfa) -- (priv);
\draw[arr] (priv) -- (ok);
\end{tikzpicture}
\end{center}
\legende{Steps to create and secure a social network account.}
\end{popfigure}

\begin{methode}[Checklist for securing a social network account]
\begin{enumerate}
\item Choose a unique password of at least 12 characters; store it safely (memory or a reputable password manager), never in a chat.
\item Activate two-factor authentication and keep your phone number up to date.
\item Run a \emph{privacy audit}: set the default audience of your posts to ``friends'', limit who can look you up by phone or e-mail, and review who can send you requests.
\item Review \emph{third-party applications} connected to your account and remove any you no longer use.
\item Check the list of active sessions/devices and log out of devices you do not recognise.
\item Repeat this audit at least once per term, and immediately after any suspicious alert.
\end{enumerate}
\end{methode}

\courssub{Building and Managing a Profile}

Your \cle{profile} is your public shop window: profile photo, cover image, biography (``bio'') and personal details. A good rule is to separate \emph{identity} information (a clear, decent photo, your name, your general interests) from \emph{sensitive} information that should stay private: your exact home address, your phone number, your daily routine, your national ID number, examination results slips, or photos of official documents. Recruiters, school authorities and strangers all look at profiles; publish only what you would accept to see on a public notice board in Yaoundé.

\courssub{Managing Connections}

Social networks grow through \cle{connections}. Depending on the platform you \emph{send or accept friend requests} (a mutual link) or you \emph{follow} accounts (a one-way subscription). Three protective tools exist everywhere:

\begin{itemize}
\item \textbf{Ignoring or declining} a request from someone you do not know;
\item \textbf{Blocking} a user, which prevents them from seeing your profile or contacting you at all;
\item \textbf{Reporting} a user or content to the platform when it is abusive, fraudulent or illegal.
\end{itemize}

\begin{attention}[Common mistake: unknown requests and public location]
Accepting every friend request feels polite but is dangerous: fake profiles collect your photos and personal details to build scams or impersonate you. Similarly, publishing your real-time location (``checked in at\ldots'') tells strangers exactly where you are and when your home is empty. Accept only people you genuinely know, and share your location, if at all, only with close friends.
\end{attention}

\courssub{Publishing Content and Choosing an Audience}

Publishing takes several forms: \cle{text posts}, \cle{photos}, \cle{videos}, and \cle{stories} (short-lived content that disappears after about 24 hours). The crucial skill is controlling the \cle{audience} of each post. Most platforms offer concentric levels of visibility, from the most private to the most public.

\begin{popfigure}
\begin{center}
\begin{tikzpicture}
\fill[popblue!8]  (0,0) circle (3.4);
\fill[popblue!18] (0,0) circle (2.5);
\fill[popblue!30] (0,0) circle (1.6);
\fill[poporange!45] (0,0) circle (0.75);
\draw[popblue, thick] (0,0) circle (3.4);
\draw[popblue, thick] (0,0) circle (2.5);
\draw[popblue, thick] (0,0) circle (1.6);
\draw[poporange, thick] (0,0) circle (0.75);
\node[font=\small\bfseries] at (0,0) {Only me};
\node[font=\small] at (0,1.15) {Friends};
\node[font=\small] at (0,2.05) {Friends of friends};
\node[font=\small] at (0,3.0) {Public};
\node[font=\footnotesize, text=popmuted, align=center] at (0,-3.95) {Each larger circle includes all smaller ones.};
\end{tikzpicture}
\end{center}
\legende{Audience levels of a post: from ``only me'' (centre) to ``public'' (outer circle).}
\end{popfigure}

\begin{exemplebox}[Choosing the right audience]
A photo of your family celebration is appropriate for \emph{friends}; a public announcement about a school club event can be \emph{public}; a personal diary-style reflection should be \emph{only me} or a custom list. Remember the direction of risk: you can always widen an audience later, but you can never ``un-publish'' content that has been copied.
\end{exemplebox}

\courssub{Interacting: Likes, Comments, Shares, Tags, Hashtags, Groups and Pages}

Interaction is the engine of social networks. You can \cle{like} (or react to) a post, \cle{comment} on it, \cle{share} it to your own audience, and \cle{tag} people with the \texttt{@} symbol so they are notified. A \cle{hashtag} (\texttt{\#}) labels a post with a topic, making it discoverable by anyone following that topic. Beyond individual posts, users can create \cle{groups} (communities of members discussing a shared interest, with membership control) and \cle{pages} (public showcases for a business, school, artist or association). The creator of a group or page becomes its \cle{administrator}: they approve members, moderate content, appoint moderators and are responsible for what happens inside.

\begin{attention}[Interactions are public acts]
Liking, commenting or sharing is not neutral: your name appears, and a share republishes content to \emph{your} audience under \emph{your} responsibility. Sharing an unverified rumour or an insulting comment can engage your liability, even if you did not write the original post.
\end{attention}

\courssub{Privacy Settings}

Beyond the audience of each post, the platform's \cle{privacy settings} let you control globally: who can see your future and past posts, who can send you friend requests or messages, whether search engines may link to your profile, whether \cle{facial recognition} may be used to suggest tags of you in photos, whether your \cle{location} is attached to posts, and which \cle{third-party applications} (games, quizzes, external sites) may read your data. Third-party apps are a frequent leak: a harmless-looking quiz may harvest your friend list and personal details. Grant applications the minimum permissions and disconnect those you no longer use.

\begin{savaistu}[Cameroon's legal framework]
Cameroon's Law No. 2010/012 on cybersecurity and cybercrime criminalises unauthorised access to personal data, identity theft and online fraud. The \cle{ANTIC} (Agence Nationale des Technologies de l'Information et de la Communication) is the regulatory body that oversees digital security and can be contacted to report cybercrimes. The principle of \emph{data minimisation} --- collecting only what is necessary --- is a cornerstone of good privacy practice worldwide.
\end{savaistu}

\courssub{Digital Identity and E-Reputation}

\begin{definition}[Digital footprint and e-reputation]
The \cle{digital footprint} is the complete set of traces a person leaves on the Internet, both \emph{active} (posts, photos, comments deliberately published) and \emph{passive} (data collected by platforms, tags by others, browsing records, metadata). The \cle{e-reputation} is the image of a person that results from this footprint --- what someone concludes about you after searching your name online.
\end{definition}

Your e-reputation has concrete consequences. University admission officers and employers increasingly consult candidates' public profiles; an old insulting post, a photo of irresponsible behaviour or hateful comments can cost you a scholarship or a job in Douala or abroad. Conversely, a clean, constructive profile (school projects, volunteering, a professional network page) becomes an asset. The \cle{right to be forgotten} is the possibility, recognised in several legal frameworks (notably the EU's GDPR), of requesting the removal or de-indexing of personal information that is outdated or harmful; in practice, prevention is far easier than deletion, because copies spread.

\begin{propriete}[Golden rule of publication]
Anything published online can be copied, screenshot and redistributed beyond your control, even if you delete it afterwards or restrict its audience. (Not examinable as a proof --- it is a practical law of digital life.)
\end{propriete}

\begin{savaistu}[Think before you post]
Before publishing, apply the three-question test: \emph{Is it true? Is it kind or useful? Would I accept my parents, my principal and a future employer seeing it in ten years?} If any answer is no, do not post. A screenshot takes one second and survives any deletion.
\end{savaistu}

\courssub{Security Threats Specific to Social Networks}

\begin{itemize}
\item \textbf{Account hacking}: attackers guess or steal passwords, then use your account to defraud your contacts. Defence: strong unique password plus 2FA.
\item \textbf{Phishing}: fake login pages or messages (``your account will be closed, click here'') that steal your credentials. Defence: never log in through a link received by message; type the address yourself.
\item \textbf{Social engineering}: manipulation techniques (pretexting, baiting, tailgating) that exploit human psychology to obtain confidential information.
\item \textbf{Identity theft and fake profiles}: someone copies your photos and name to create a duplicate account and ask your friends for money. Defence: keep your audience restricted, and report impostor accounts immediately to the platform and, if necessary, to ANTIC.
\item \textbf{Scams}: fraudulent investment offers, fake job adverts, and the \emph{money-doubling} frauds circulating on Cameroonian WhatsApp and Facebook groups, where victims are promised that a mobile money transfer will be ``multiplied''. No legitimate system doubles money; anyone asking you to send money first in order to receive more is a fraudster.
\end{itemize}

\begin{attention}[Never share passwords, OTP codes or ID photos]
No legitimate service --- not the platform, not your bank, not MTN or Orange --- will ever ask for your password or your one-time 2FA code in a chat or by phone. A ``friend'' (often a hacked account) who asks you to forward an SMS code is stealing your account. Likewise, never send photos of your national ID card, voter's card or school documents in a messaging conversation: they are raw material for identity theft.
\end{attention}

\begin{exoresolu}[Analysing a suspicious message]
Amina, an Upper Sixth student in Limbe, receives a Facebook message from the account of her classmate Beri: ``Please, I sent you a code by mistake, forward me the SMS you just received.'' At the same moment, Amina receives an SMS containing a 6-digit code from the social network. What is happening, and what should Amina do?
\tcblower
\textbf{Analysis.} Beri's account has almost certainly been hacked. The attacker is trying to take over \emph{Amina's} account: they triggered a login or password reset on Amina's account, which caused the network to send the 2FA code to Amina's phone. The code is the last barrier; forwarding it hands the account to the attacker.
\textbf{Correct actions:}
\begin{enumerate}
\item Do \emph{not} forward the code to anyone --- a genuine service never asks for it.
\item Contact Beri by another channel (a phone call, in person) to warn him that his account is compromised.
\item Report the suspicious conversation to the platform and block the sender if necessary.
\item Check her own account: active sessions, password and 2FA still enabled.
\end{enumerate}
\textbf{Lesson:} the OTP code is as secret as the password itself.
\end{exoresolu}

\begin{exercice}[Securing a new account]
Claude has just created his first social network account to promote his school's science club. (a) List four steps he should take immediately to secure the account. (b) The club wants maximum visibility: what audience should Claude choose for the club's announcements, and what audience for photos of club members? Justify. (c) A member wants to post the club's mobile money number publicly ``for donations''. Give two risks of this idea. (d) If Claude suspects the account has been compromised, which Cameroonian authority can he contact?
\end{exercice}

\begin{corrige}[Exercise 1]
(a) Choose a strong unique password; activate 2FA; verify the e-mail or phone number; set the default audience and run a privacy audit (who can contact him, app permissions, active sessions).
(b) Club announcements: \emph{public}, so the club gains visibility. Photos of members: \emph{friends} or a custom audience, and only with the members' explicit consent, because identifiable images of minors or students should not be exposed to everyone.
(c) Publishing a mobile money number publicly exposes the club to scam calls and phishing attempts targeting that number, and allows fraudsters to impersonate the club and collect fake ``donations'' in its name. A controlled channel (official page with verified contact, or direct contact with the treasurer) is safer.
(d) He can report the incident to \cle{ANTIC} (Agence Nationale des Technologies de l'Information et de la Communication), which handles cybersecurity incidents and cybercrime in Cameroon.
\end{corrige}

\begin{aretenir}[Key points of the section]
\begin{itemize}
\item Secure every account: strong unique password, verified contact address, and two-factor authentication.
\item Control the audience of every post: only me, friends, friends of friends, or public.
\item Publish only what you accept to see copied forever; your digital footprint builds your e-reputation, which affects studies and employment.
\item Block and report fake profiles; never share passwords, OTP codes or photos of official documents.
\item Treat ``too good to be true'' offers --- especially money-doubling schemes --- as frauds.
\item Know your rights and the legal framework: Law No. 2010/012, ANTIC, data minimisation, right to be forgotten.
\end{itemize}
\end{aretenir}

\coursec{Collaborative Environments and Tools for Teamwork}

In the previous section we learnt how to create an account on a social network, configure privacy settings and communicate safely. Social networks connect people around \emph{conversations}; but very often, at school, at university or at work, people do not simply want to \emph{talk} online --- they want to \emph{produce something together}: a group assignment, a budget spreadsheet, a project report, a revision document. This is where \cle{collaborative environments} come in. They are the natural continuation of social networking: from \emph{connecting} people we move to making people \emph{work together}.

\courssub{What is a Collaborative Environment?}

\textbf{Intuition first.} Imagine three classmates in Yaound\'e, Bamenda and Buea who must submit a joint ICT project. Twenty years ago, they would have had to meet physically, or exchange handwritten pages. Today, they open the \emph{same} document on their phones, each writes a part, and everyone sees the others' changes appear live. The document exists only once, on a server ``in the cloud'', and it is always up to date. The digital space that makes this possible is a collaborative environment.

\begin{definition}[Collaborative environment (groupware)]
A \cle{collaborative environment} (also called \cle{groupware}) is a digital space --- a set of software tools and services, usually hosted on remote servers (the \cle{cloud}) --- that allows \textbf{several users to work together on common tasks}, whether they are in the same room or thousands of kilometres apart, and whether they work \textbf{at the same time (synchronously)} or \textbf{at different times (asynchronously)}.
\end{definition}

\begin{aretenir}[Key features of a collaborative environment]
A good collaborative environment provides: a \textbf{shared workspace} (common files, one single version), \textbf{communication tools} (chat, comments, videoconference), \textbf{coordination tools} (tasks, calendars, deadlines) and a \textbf{history} that records who did what and when (\cle{traceability}).
\end{aretenir}

\courssub{Categories of Collaborative Tools}

Collaborative tools can be grouped into families according to their main purpose. The mind map below summarises the principal categories, with two examples each.

\begin{popfigure}
\begin{tikzpicture}[grow cyclic, mindmap, concept color=popblue,
  level 1 concept/.append style={level distance=4.2cm, sibling angle=72},
  every node/.style={concept, text=white, font=\small}]
\node[font=\normalsize\bfseries, align=center]{Collaborative tools}
  child[concept color=popgreen] { node {Office suites in the cloud}
      child[concept color=popgreen!70] { node[font=\scriptsize] {Google Workspace} }
      child[concept color=popgreen!70] { node[font=\scriptsize] {Microsoft 365} } }
  child[concept color=poporange] { node {Cloud storage \& sharing}
      child[concept color=poporange!70] { node[font=\scriptsize] {Google Drive} }
      child[concept color=poporange!70] { node[font=\scriptsize] {Dropbox} } }
  child[concept color=poppurple] { node {Wikis}
      child[concept color=poppurple!70] { node[font=\scriptsize] {Wikipedia} }
      child[concept color=poppurple!70] { node[font=\scriptsize] {Wikibooks} } }
  child[concept color=poppink] { node {Project management}
      child[concept color=poppink!70] { node[font=\scriptsize] {Trello} }
      child[concept color=poppink!70] { node[font=\scriptsize] {Asana} } }
  child[concept color=popgold] { node {Videoconferencing}
      child[concept color=popgold!70] { node[font=\scriptsize] {Zoom} }
      child[concept color=popgold!70] { node[font=\scriptsize] {Google Meet} } };
\end{tikzpicture}
\legende{Mind map of the main categories of collaborative tools, with two examples each.}
\end{popfigure}

\begin{itemize}
\item \textbf{Office suites in the cloud} (Google Workspace, Microsoft 365): create and edit text documents, spreadsheets and presentations directly in a browser, with several people at once.
\item \textbf{Cloud storage and file sharing} (Google Drive, OneDrive, Dropbox): store files online and share them through links, so that everyone accesses the same version from any device.
\item \textbf{Wikis} (Wikipedia): websites whose pages can be created and modified by any authorised member, with a complete history of every change.
\item \textbf{Project management tools} (Trello, Asana): organise work into tasks and boards, assign responsibilities and track progress.
\item \textbf{Videoconferencing} (Zoom, Google Meet, Microsoft Teams): hold live meetings with sound, video and screen sharing.
\end{itemize}

\begin{savaistu}[Wikis and Wikipedia]
The word \emph{wiki} comes from a Hawaiian expression meaning ``quick''. Wikipedia, the most famous wiki, is written entirely by volunteers: any registered contributor can improve an article, and every modification is recorded in the page history, so vandalism or errors can be undone in seconds. The same principle is used by companies to build internal knowledge bases.
\end{savaistu}

\courssub{Collaborative Editing: One Document, Many Authors}

The heart of a collaborative environment is \cle{collaborative editing}: several users work on \textbf{one single shared document} instead of exchanging copies. Four mechanisms make this possible:

\begin{itemize}
\item \textbf{Simultaneous editing.} Each connected user sees the cursors and changes of the others in real time. The server synchronises all modifications so that everyone always sees the latest version.
\item \textbf{Version history.} The platform automatically records every state of the document. You can see who changed what, and restore any previous version --- a safety net against mistakes.
\item \textbf{Comments and suggestions.} Instead of modifying the text directly, a collaborator can attach a comment to a passage or propose a change in ``suggestion mode'', which the owner can accept or reject.
\item \textbf{Sharing permissions.} The owner controls what each person may do: \textbf{view} (read only), \textbf{comment} (read and annotate) or \textbf{edit} (full modification).
\end{itemize}

\begin{popfigure}
\begin{tikzpicture}[font=\small]
% server cloud
\node[cloud, draw=popblue, fill=popblueL, cloud puffs=12, cloud puff arc=110,
      minimum width=4.6cm, minimum height=2.6cm, align=center] (srv) at (0,0)
      {\textbf{Cloud server}\\[1pt] \scriptsize shared document\\ \scriptsize version history};
% students
\node[draw=popgreen, fill=popgreenL, rounded corners, align=center, minimum width=2.6cm] (a) at (-5,2.6) {\textbf{Amina}\\ \scriptsize Yaound\'e};
\node[draw=poporange, fill=white, rounded corners, align=center, minimum width=2.6cm] (b) at (0,3.6) {\textbf{Bertrand}\\ \scriptsize Bamenda};
\node[draw=poppurple, fill=white, rounded corners, align=center, minimum width=2.6cm] (c) at (5,2.6) {\textbf{Clarisse}\\ \scriptsize Buea};
% arrows
\draw[<->, very thick, popgreen] (a.south east) -- (srv.north west) node[midway, above left, popgreen, font=\scriptsize, align=center]{writes \&\\ receives};
\draw[<->, very thick, poporange] (b.south) -- (srv.north) node[midway, right, poporange, font=\scriptsize, align=left]{sync};
\draw[<->, very thick, poppurple] (c.south west) -- (srv.north east) node[midway, above right, poppurple, font=\scriptsize, align=center]{writes \&\\ receives};
\end{tikzpicture}
\legende{Three students edit the same cloud document simultaneously; the server synchronises every change so that all see one single, up-to-date version.}
\end{popfigure}

\begin{methode}[Sharing a document with the correct permissions and tracking versions]
\begin{enumerate}
\item Create the document in a cloud suite (e.g.\ Google Docs) --- never as an attachment to an e-mail.
\item Click \textbf{Share} and add your collaborators by e-mail address, or generate a sharing link.
\item For each person, choose the right permission: \textbf{Viewer} for someone who must only read, \textbf{Commenter} for a proofreader, \textbf{Editor} for a co-author.
\item Avoid ``anyone with the link can edit'' for sensitive work: restrict editing to named accounts.
\item While working, use \textbf{suggestion mode} or comments for proposed changes instead of overwriting others' text.
\item Regularly open \textbf{File $\rightarrow$ Version history} to check contributions, name important versions, and restore an earlier state if something was deleted by mistake.
\end{enumerate}
\end{methode}

\begin{attention}[Version chaos: the e-mail attachment trap]
A very common mistake is to write a document, send it \textbf{as an e-mail attachment} to five teammates, and ask each one to complete it. You then receive five \emph{different} files named ``project-final'', ``project-final2'', ``project-REAL-final''\dots{} Nobody knows which is the latest version, and merging the changes wastes hours. \textbf{Rule: share one single cloud file; never circulate copies.} One document, one link, one truth.
\end{attention}

\courssub{Synchronous versus Asynchronous Collaboration}

\begin{definition}[Synchronous and asynchronous collaboration]
\textbf{Synchronous} collaboration happens when all participants work or communicate \textbf{at the same time}: a Zoom lesson, a live chat, simultaneous editing of a document. \textbf{Asynchronous} collaboration happens when participants contribute \textbf{at different times}: e-mail, forum messages, comments left in a document, tasks updated on a Trello board --- each person reads and replies when available.
\end{definition}

\begin{popfigure}
\begin{tikzpicture}[font=\small]
\node[draw=popblue, fill=popblueL, rounded corners, align=center, text width=6.4cm, minimum height=1.1cm] (s) at (-3.6,3.2) {\textbf{Synchronous}\\ \scriptsize same time, immediate feedback};
\node[draw=poporange, fill=white, rounded corners, align=center, text width=6.4cm, minimum height=1.1cm] (a) at (3.6,3.2) {\textbf{Asynchronous}\\ \scriptsize different times, flexible};
% examples
\node[align=left, text width=6.4cm] at (-3.6,1.2) {\scriptsize $\bullet$ Zoom / Meet lesson\\ \scriptsize $\bullet$ live chat, phone call\\ \scriptsize $\bullet$ simultaneous editing\\ \scriptsize $\bullet$ screen-sharing session};
\node[align=left, text width=6.4cm] at (3.6,1.2) {\scriptsize $\bullet$ e-mail, forum posts\\ \scriptsize $\bullet$ document comments\\ \scriptsize $\bullet$ Trello task updates\\ \scriptsize $\bullet$ recorded video lessons};
% when to use
\node[draw=popgreen, fill=popgreenL, rounded corners, align=center, text width=6.4cm] at (-3.6,-1.2) {\scriptsize \textbf{Use when:} urgent decision, brainstorming, teaching, resolving a conflict quickly};
\node[draw=poppurple, fill=white, rounded corners, align=center, text width=6.4cm] at (3.6,-1.2) {\scriptsize \textbf{Use when:} different timetables, limited connectivity, thoughtful written feedback};
\end{tikzpicture}
\legende{Synchronous versus asynchronous collaboration: examples and typical situations of use.}
\end{popfigure}

Neither mode is ``better'': effective teams combine both. A group might meet synchronously on Google Meet to plan the work, then collaborate asynchronously in the shared document during the week --- which is precisely how most Cameroonian students organise group assignments around their class timetables and connectivity constraints.

\courssub{E-learning Platforms as Collaborative Environments}

E-learning platforms are complete collaborative environments dedicated to education. \textbf{Moodle}, used by many Cameroonian universities, lets a teacher publish courses, collect assignments, run forums and quizzes, while students submit work and discuss. \textbf{Google Classroom} organises classes, distributes and grades assignments, and links directly to Google Docs. Even \textbf{WhatsApp and Telegram study groups} --- extremely widespread in Cameroonian schools and universities --- act as lightweight collaborative environments: teachers share documents and announcements, students ask questions, and voice notes replace the classroom when needed. During periods of school disruption, these tools have allowed teaching to continue at a distance.

\courssub{Advantages, Limits and Challenges}

\begin{center}
\begin{tabularx}{\linewidth}{|l|X|}\hline
\rowcolor{popblueL} \textbf{Advantages} & \textbf{Explanation} \\ \hline
Pooling of skills & Each member contributes his or her strength (writing, calculation, design). \\ \hline
Distance teamwork & Members in different towns work as if in the same room. \\ \hline
Time saving & No travelling, no merging of copies; work progresses in parallel. \\ \hline
Traceability & Version history shows who contributed what --- useful for fair assessment. \\ \hline
Continuity of work & Work goes on despite disruptions (travel, illness, school interruptions). \\ \hline
\end{tabularx}
\end{center}

\begin{attention}[Limits and challenges in the Cameroonian context]
Collaborative work online requires a \textbf{reliable Internet connection}, and mobile data remains a significant cost for many families. The \textbf{digital divide} means some teammates may have no computer or only irregular access. Other challenges include \textbf{conflicts over contributions} (a member who does nothing, or disagreements about content) and \textbf{dependence on the platform}: if the service closes or your account is blocked, your workspace disappears --- hence the importance of regular local backups.
\end{attention}

\courssub{Organising Effective Group Work Online}

Tools alone do not guarantee success; organisation does. An effective online group: \textbf{defines roles} (a coordinator, a secretary, a proofreader\dots), \textbf{sets clear rules} (which tool for what, response times, netiquette), \textbf{fixes deadlines} on a shared calendar or Trello board, and works in \textbf{one shared workspace} rather than scattered files. The coordinator checks the version history to ensure everyone contributes.

\begin{exoresolu}[Choosing tools and permissions for a group project]
Four Upper Sixth students --- two in Douala, one in Yaound\'e, one in Garoua --- must produce a 10-page report and a spreadsheet of results in three weeks. One member has only a phone and limited data. Propose a collaborative organisation: tools, permissions, and working mode.
\tcblower
\textbf{Proposed solution.}
\begin{itemize}
\item \textbf{Workspace:} one shared folder on Google Drive containing \emph{one} Google Doc (report) and \emph{one} Google Sheet (results) --- never e-mail attachments.
\item \textbf{Permissions:} the three writers get \emph{edit} rights on the Doc; the member with limited data gets \emph{comment} rights so he can review without heavy editing sessions; the teacher receives a \emph{view-only} link.
\item \textbf{Mode:} one short \textbf{synchronous} Google Meet session per week (30 minutes, cameras off to save data) to decide tasks; the rest is \textbf{asynchronous}: writing in the shared document, comments, and a Trello board listing tasks with owners and deadlines.
\item \textbf{Roles:} a coordinator (checks version history weekly), a secretary (meeting notes), a proofreader (suggestion mode), a data manager (spreadsheet).
\item \textbf{Safety:} the coordinator downloads a backup copy each week and names the final version in the version history before submission.
\end{itemize}
This organisation respects the connectivity constraint, guarantees a single up-to-date version, and makes every contribution traceable.
\end{exoresolu}

\begin{exercice}[Analysing a collaborative situation]
A class WhatsApp group is used to prepare a science fair. Files are sent as attachments, two students argue about who wrote a paragraph, and the final poster is lost when a phone is stolen. Identify \textbf{four} weaknesses of this organisation and, for each, propose a remedy based on the collaborative tools studied in this section.
\end{exercice}

\begin{corrige}[Exercise 1]
\begin{enumerate}
\item \textbf{Files sent as attachments} $\rightarrow$ version chaos. Remedy: create the documents in Google Docs/Slides and share one link with appropriate permissions.
\item \textbf{Dispute over authorship} $\rightarrow$ no traceability. Remedy: work in the shared document so the version history records who wrote what.
\item \textbf{Poster lost with the phone} $\rightarrow$ no backup, dependence on one device. Remedy: store the file in cloud storage (Drive/Dropbox) so it survives the loss of any device.
\item \textbf{Everything mixed in one chat} $\rightarrow$ no task tracking. Remedy: add a Trello board (or a simple shared checklist) with tasks, responsible members and deadlines, and keep WhatsApp only for quick synchronous announcements.
\end{enumerate}
\end{corrige}

\begin{aretenir}[To remember]
A collaborative environment lets several users work together on common tasks, synchronously or asynchronously, around \textbf{one shared version} of each file. Master the categories of tools, the three permission levels (view / comment / edit), the version history, and the discipline of good organisation (roles, rules, deadlines, shared workspace) --- and never confuse a chat group with a true collaborative workspace.
\end{aretenir}

\coursec{Risks, Ethics, Netiquette and Legal Framework}

In the previous section we saw how collaborative environments --- shared documents, cloud workspaces, group chats --- allow teams to work together efficiently. But the same tools that connect us also expose us: a message can be forwarded to thousands of people in seconds, a photo can be copied forever, and a careless comment can become a legal offence. This final section therefore asks the essential question: \emph{how do we use social networks and collaborative tools safely, ethically and legally?} We study the advantages and dangers, the rules of good behaviour online (netiquette), intellectual property, and the Cameroonian legal framework, before concluding with responsible digital citizenship.

\courssub{Advantages and dangers of social networks: a balanced view}

\begin{experience}[Observation: one tool, two faces]
A student in Bamenda uses a WhatsApp study group to share past GCE questions and explain difficult exercises to classmates --- a clear benefit. The same week, a rumour spread in another WhatsApp group causes panic in her neighbourhood, and a classmate is mocked cruelly on Facebook because of her appearance. The \emph{same technology} produced learning, panic and harm. The difference lies not in the tool but in how it is used.
\end{experience}

Social networks offer real, well-documented advantages:

\begin{itemize}
\item \textbf{Communication:} keeping in touch with family and friends across Cameroon and the diaspora at very low cost.
\item \textbf{Information:} following news, public announcements (for example from ministries or councils) and events in real time.
\item \textbf{Learning:} study groups, tutorial videos, online courses and professional communities.
\item \textbf{Business promotion:} a tailor in Douala or a farmer in Bafoussam can advertise products on Facebook or WhatsApp Business and reach customers without paying for a shop front.
\item \textbf{Civic participation:} citizens can debate public issues, report problems and take part in community life.
\item \textbf{Job hunting:} professional networks such as LinkedIn, and job offers circulated in groups, connect employers and job seekers.
\end{itemize}

Against these advantages stand serious \cle{dangers and abuses}:

\begin{itemize}
\item \textbf{Cyberbullying and online harassment:} repeated mockery, threats or humiliation of a person online.
\item \textbf{Hate speech:} messages attacking people because of their ethnic group, religion, region or gender --- a particularly sensitive danger in a diverse country like Cameroon.
\item \textbf{Fake news and disinformation:} false stories created to deceive, manipulate opinion or earn money from clicks.
\item \textbf{Scams and fraud:} fake job offers, fake mobile-money promotions, phishing messages pretending to come from MTN, Orange or a bank.
\item \textbf{Addiction and time waste:} endless scrolling that eats study time and sleep.
\item \textbf{Exposure to inappropriate content:} violent or adult content, especially harmful to younger users.
\item \textbf{Collaborative environment risks:} accidental sharing of sensitive files, unauthorised edits in shared documents, version conflicts, and excessive permissions granted to team members.
\end{itemize}

\begin{definition}[Cyberbullying]
\cle{Cyberbullying} is the repeated, intentional use of digital technologies (social networks, messaging apps, forums) to threaten, humiliate, insult or exclude a person. Unlike face-to-face bullying, it can happen at any hour, reach a huge audience, and leave permanent traces (screenshots, shared photos), which makes it especially harmful to victims.
\end{definition}

\begin{definition}[Fake news]
\cle{Fake news} (disinformation) is false or misleading information presented as genuine news, created and spread deliberately to deceive the public, damage a reputation, influence opinions or generate advertising revenue. It is distinct from a simple honest mistake: fake news involves the \emph{intention} to deceive or a careless disregard for the truth.
\end{definition}

\begin{popfigure}
\begin{tikzpicture}[scale=0.95, every node/.style={transform shape}]
\node[draw=popgreen, fill=popgreenL, rounded corners, text width=6.2cm, align=left, inner sep=8pt] (adv) at (0,0) {%
\textbf{\textcolor{popdark}{ADVANTAGES}}\\[2pt]
$\bullet$ Cheap communication\\
$\bullet$ Fast access to information\\
$\bullet$ Learning and study groups\\
$\bullet$ Business promotion\\
$\bullet$ Civic participation\\
$\bullet$ Job hunting};
\node[draw=poppink, fill=poppink!12, rounded corners, text width=6.2cm, align=left, inner sep=8pt] (dan) at (9.2,0) {%
\textbf{\textcolor{popdark}{DANGERS}}\\[2pt]
$\bullet$ Cyberbullying, harassment\\
$\bullet$ Hate speech\\
$\bullet$ Fake news, disinformation\\
$\bullet$ Scams and fraud\\
$\bullet$ Addiction, time waste\\
$\bullet$ Inappropriate content\\
$\bullet$ Collaborative tool misuse};
\draw[line width=2.5pt, popdark] (4.6,-2.6) -- (4.6,2.6);
\draw[line width=2pt, popdark] (1.2,2.0) -- (8.0,2.0);
\fill[popdark] (4.6,2.0) circle (3pt);
\draw[line width=1.5pt, poporange] (4.6,-2.6) -- (3.6,-3.4) -- (5.6,-3.4) -- cycle;
\node[text width=3.4cm, align=center, popmuted] at (4.6,-4.1) {The user tips the balance};
\end{tikzpicture}
\legende{Social networks and collaborative tools: the same platforms carry advantages and dangers; responsible use tips the balance.}
\end{popfigure}

\courssub{Privacy issues: your data is the product}

Why are Facebook, TikTok or Instagram free? Because the real business model is the \emph{exploitation of personal data}. When you use a social network, the platform collects:

\begin{itemize}
\item the information you give directly (name, age, school, photos, posts, likes);
\item your behaviour (what you watch, for how long, what you click, who you talk to);
\item technical data (device type, IP address) and often your \cle{geolocation}.
\end{itemize}

These data are used to build a detailed \cle{profile} of you --- your interests, habits, opinions, even your mood --- which is then used for \cle{targeted advertising} and sometimes sold or shared with third parties. Geolocation tracking can reveal where you live, where you study and where you go every day. In collaborative environments, shared folders and documents may expose metadata (author names, edit history, timestamps) that reveals more than intended.

\begin{attention}[If it is free, you are the product]
Never publish your home address, your daily routine, your exam slip with personal identifiers, or intimate photos. Once published, you lose control of the data: copies can be made, stored and resold without your consent. Adjust the privacy settings of every account and review them regularly. In shared drives, set permissions to \emph{view only} unless editing is strictly necessary, and audit access lists periodically.
\end{attention}

\courssub{Fake news: mechanisms of virality and verification}

Fake news spreads fast because it is designed to trigger \emph{emotion} --- fear, anger, surprise --- and emotional content is shared without reflection. Social network algorithms amplify whatever gets reactions, so a shocking falsehood often travels further than a careful truth. Closed messaging groups (WhatsApp, Telegram) make verification harder because messages arrive from people we trust.

\begin{methode}[The SIFT technique to verify information]
Before believing or sharing any surprising piece of news, apply the four moves of the \cle{SIFT} method (developed by digital literacy expert Mike Caulfield):
\begin{enumerate}
\item \textbf{Stop:} do not share immediately. Ask yourself: do I actually know this is true? Am I reacting emotionally?
\item \textbf{Investigate the source:} who published this? Is it a known, credible media outlet or institution, or an anonymous page? Check the ``About'' section of the site or account; look for a verifiable physical address and editorial team.
\item \textbf{Find better coverage:} search the same story on other independent, trusted sources. If only one obscure page reports it, be suspicious. Use fact-checking sites where available (e.g., AFP Fact Check, Dubawa).
\item \textbf{Trace claims:} go back to the original context. Use a \emph{reverse image search} (for example Google Images, TinEye) to check whether a photo is old or taken elsewhere; check the date; find the original quote or document.
\end{enumerate}
If any step fails, \textbf{do not share}.
\end{methode}

\begin{popfigure}
\begin{tikzpicture}[node distance=0.55cm,
box/.style={draw=popblue, fill=popblueL, rounded corners, text width=4.6cm, align=center, inner sep=6pt},
dec/.style={draw=poporange, fill=popgold!20, diamond, aspect=2.2, text width=3.2cm, align=center, inner sep=2pt},
res/.style={draw=popgreen, fill=popgreenL, rounded corners, text width=3.4cm, align=center, inner sep=6pt},
bad/.style={draw=poppink, fill=poppink!15, rounded corners, text width=3.4cm, align=center, inner sep=6pt},
arr/.style={-{Stealth}, thick, popdark}]
\node[box] (see) {You receive surprising news};
\node[dec, below=of see] (src) {Source credible and identified?};
\node[dec, below=of src] (date) {Date and context correct?};
\node[dec, below=of date] (cross) {Confirmed by other trusted sources?};
\node[res, below=of cross] (share) {Share responsibly, citing the source};
\node[bad, right=2.6cm of date] (noshare) {Do NOT share; report if dangerous};
\draw[arr] (see) -- (src);
\draw[arr] (src) -- node[right, popmuted]{yes} (date);
\draw[arr] (date) -- node[right, popmuted]{yes} (cross);
\draw[arr] (cross) -- node[right, popmuted]{yes} (share);
\draw[arr] (src.east) -- ++(1.1,0) node[above, popmuted, pos=0.6]{no} |- (noshare.north);
\draw[arr] (date.east) -- node[above, popmuted]{no} (noshare.west);
\draw[arr] (cross.east) -- ++(1.1,0) node[above, popmuted, pos=0.6]{no} |- (noshare.south);
\end{tikzpicture}
\legende{Verification flowchart: check the source, the date and independent coverage before sharing anything.}
\end{popfigure}

\begin{attention}[Forwarding is publishing]
A very common mistake is to think: ``I only forwarded the message, I did not write it, so I am not responsible.'' This is false. In Cameroon, forwarding unverified or false information in WhatsApp groups can be punished by law, because forwarding contributes to the \emph{publication} of false news. Group administrators can also be held responsible for content they knowingly allow. When in doubt, do not forward.
\end{attention}

\courssub{Netiquette: the rules of polite behaviour online}

\begin{definition}[Netiquette and its five golden rules]
The \cle{netiquette} (contraction of \emph{network} and \emph{etiquette}) is the set of rules of polite, respectful and responsible behaviour on the Internet. Its five golden rules are:
\begin{enumerate}
\item \textbf{Respect others:} no insults, no mockery, no hate speech; respect opinions different from yours.
\item \textbf{Do not shout:} writing in ALL CAPITAL LETTERS is perceived as shouting and aggression.
\item \textbf{Do not spam:} never send unwanted advertisements, chain messages or repeated messages.
\item \textbf{Respect privacy:} never publish someone's photo, phone number or private conversation without their consent.
\item \textbf{Cite your sources:} when you reuse a text, image or idea, mention the author; do not present others' work as your own.
\end{enumerate}
\end{definition}

\courssub{Intellectual property online}

Everything published online --- photos, music, videos, texts, software --- is protected by \cle{copyright}, which belongs by default to its creator. The fact that an image is easy to download does \emph{not} make it free to use. Sharing a song, reposting a photographer's picture, or copying a text into your assignment without permission and without citing the source is a violation of intellectual property.

In collaborative work, a specific danger is \cle{plagiarism}: presenting someone else's words or ideas as your own. When a group produces a shared document, each member must cite borrowed material, and the group must never simply copy-paste from websites. Teachers and examiners can detect plagiarism easily with search engines and specialised software.

\begin{exemplebox}[Citing correctly in a shared document]
Your group uses a map of Cameroon's fibre-optic backbone found on the website of a telecommunications operator. Correct practice: insert the image, then write beneath it ``Source: [name of operator] website, consulted on [date]'', and ask for permission if the work will be published. Simply copying the image with no mention is plagiarism.
\end{exemplebox}

\courssub{The legal framework in Cameroon}

\begin{aretenir}[Law No. 2010/012 on cybersecurity and cybercriminality]
In Cameroon, \cle{Law No. 2010/012 of 21 December 2010} relating to cybersecurity and cybercriminality defines and punishes offences committed through information and communication technologies. Among the offences it sanctions are:
\begin{itemize}
\item \textbf{Identity theft:} using another person's name, photos or accounts to deceive others;
\item \textbf{Online defamation:} publishing false statements that damage a person's honour or reputation;
\item \textbf{Publication of false news:} spreading, including by simple forwarding, information known to be false;
\item \textbf{Cyber harassment, intrusion into systems, and electronic fraud.}
\end{itemize}
Penalties include \textbf{prison sentences and heavy fines}, which can be doubled in aggravated circumstances. Online actions therefore have real, serious legal consequences.
\end{aretenir}

\begin{savaistu}[ANTIC: the guardian of cyberspace in Cameroon]
Cameroon created the National Agency for Information and Communication Technologies (\cle{ANTIC}) to, among other missions, promote the security of information systems, monitor cyber threats and support the fight against cybercrime. Victims of online offences can report them to the competent authorities, and specialised units of the police and gendarmerie handle cybercrime complaints.
\end{savaistu}

\begin{exoresolu}[Analysing a real-life situation]
\textbf{Statement.} A Form Five student creates a fake Facebook account using the name and photos of a classmate, and posts insulting messages and a false rumour about a teacher. The rumour is forwarded by several students in a WhatsApp group. Identify all the wrongdoings and the rules or laws broken.\tcblower
\textbf{Solution.}
\begin{enumerate}
\item \textbf{Fake account with a classmate's identity:} this is \emph{identity theft}, an offence under Law No. 2010/012, and a violation of the classmate's privacy (netiquette rule 4).
\item \textbf{Insulting messages:} this is \emph{cyberbullying / online harassment} and breaks the netiquette rule of respect.
\item \textbf{False rumour about the teacher:} this is \emph{online defamation} and \emph{publication of false news}, both punishable by Law No. 2010/012.
\item \textbf{Students who forwarded the rumour:} by forwarding unverified false information, they contributed to the publication of false news and can also be punished; they should have applied the SIFT method and refused to share.
\item \textbf{Correct behaviour:} the victims should keep screenshots as evidence, report the account to the platform, inform the school authorities and, if necessary, file a complaint with the police or gendarmerie.
\end{enumerate}
\end{exoresolu}

\courssub{Responsible digital citizenship}

Being a good \cle{digital citizen} means using networks positively and protecting the community:

\begin{itemize}
\item \textbf{Positive use:} share useful, verified content; encourage others; use networks to learn, create and promote honest activities.
\item \textbf{Report abuse:} use the ``report'' tools of platforms, keep evidence (screenshots), and alert adults or authorities when someone is harassed or threatened.
\item \textbf{Protect younger users:} guide younger brothers, sisters and pupils; help them configure privacy settings; warn them about strangers online and inappropriate content.
\item \textbf{Balance online and offline life:} set time limits, switch off notifications while studying, and keep real-life friendships, sport and family time.
\item \textbf{Collaborative responsibility:} in shared workspaces, respect access rights, do not delete others' contributions without agreement, and communicate changes clearly.
\end{itemize}

\begin{propriete}[The ``grandmother test'']
Before posting anything, apply the \cle{grandmother test}: \emph{if you would be ashamed to show this post to your grandmother, your family or a future employer, do not post it.} Nothing published online ever truly disappears: screenshots and archives survive even after deletion, and employers routinely check candidates' profiles.
\end{propriete}

\begin{exercice}[Applying the rules]
For each situation, state the rule or law concerned and the correct behaviour:
\begin{enumerate}
\item A friend sends you a shocking video claiming a disaster happened this morning in Yaound\'e; no media reports it.
\item You want to use a professional photographer's picture in your class blog.
\item A classmate is mocked daily in the class WhatsApp group.
\item Your team shares a Google Drive folder for a project; one member accidentally deletes a critical file.
\end{enumerate}
\end{exercice}

\begin{corrige}[Exercise 1]
\begin{enumerate}
\item Apply SIFT: stop, check the source, look for coverage by trusted media. Since nothing confirms it, do not forward; forwarding false news is punishable under Law No. 2010/012.
\item The photo is protected by copyright: ask the photographer's permission, or use a free-licence image (e.g., Creative Commons), and always cite the source.
\item This is cyberbullying: do not participate, support the victim privately, keep evidence, report to the group administrator and to a teacher or school authority.
\item Use the version history (``Manage versions'' in Google Drive) to restore the deleted file; establish a team rule to avoid direct deletion and use ``suggesting'' mode for edits.
\end{enumerate}
\end{corrige}

\begin{aretenir}[Key points of the section]
\begin{itemize}
\item Social networks bring communication, learning, business and civic benefits, but also cyberbullying, hate speech, fake news, scams and addiction.
\item Collaborative environments add risks: accidental data exposure, permission misuse, version conflicts.
\item Platforms collect, profile and monetise your personal data: protect your privacy and your geolocation.
\item Verify before sharing with the SIFT method; forwarding false information is legally punishable.
\item Follow the netiquette: respect, no ALL-CAPS, no spam, respect privacy, cite sources.
\item Copyright protects online content; plagiarism in collaborative work is a serious fault.
\item Law No. 2010/012 of 21 December 2010 punishes identity theft, online defamation and the publication of false news in Cameroon.
\item Be a responsible digital citizen: use networks positively, report abuse, protect the youngest, and apply the grandmother test before posting.
\end{itemize}
\end{aretenir}

\newpage

\coursec{Integration activity}

\courssub{Aims of the activity}

This integration activity brings together everything you have learnt in this chapter: the concepts and types of \cle{social networks}, the responsible use of accounts and \cle{privacy} settings, the choice and configuration of \cle{collaborative environments}, and the rules of \cle{netiquette}, ethics and the Cameroonian legal framework (in particular Law No. 2010/012 of 21 December 2010 on cybersecurity and cybercriminality).

By the end of the activity, you should be able to:
\begin{itemize}
\item \textbf{Choose} a collaborative platform and \textbf{justify} the choice using objective criteria (cost, connectivity, features, supervision, data protection);
\item \textbf{Create and configure} a shared digital workspace with correct folder structure, roles and \cle{sharing permissions} (view, comment, edit, owner);
\item \textbf{Co-edit} a document in real time, using comments, suggestions and \cle{version history};
\item \textbf{Draft} a charter of netiquette and security adapted to a class group;
\item \textbf{Analyse} advantages and risks of online collaboration and \textbf{relate} them to the principles of Law No. 2010/012;
\item \textbf{Present and defend} your work orally, as expected in the GCE Advanced Level practical and project assessments.
\end{itemize}

\begin{aretenir}[Skills mobilised]
Analysing a need, selecting a tool, configuring access rights, collaborating online, communicating respectfully, protecting personal data, arguing a choice, working in a team under a deadline.
\end{aretenir}

\courssub{Situation}

\textbf{Government High School Bafoussam — Upper Sixth Science, 38 students.}\par
It is the end of the First Term. In six weeks, your class will sit the \emph{mock GCE Advanced Level ICT examination}. During the last class meeting, the students complained about their revision habits: everyone photocopies different notes, some summaries circulate as photos in a WhatsApp group where they are quickly buried under hundreds of unrelated messages, two students lost their exercise books, and last term a wrong ``leaked paper'' (in fact an old paper from another school) caused panic the night before the test.

Your ICT teacher, Mrs.\ Ngum, proposes a solution: \textbf{the class will build a Collaborative Revision Hub} — a supervised digital workspace where every chapter of the syllabus has one well-organised, verified summary sheet, co-written and corrected by the students themselves.

The class is divided into \textbf{groups of 4 or 5 students}. Your group is in charge of the chapter \emph{``Using social networks and collaborative environments''}, and must also act as the pilot group whose choices will serve as a model for the others.

\textbf{Available resources and constraints (data of the situation):}
\begin{itemize}
\item The school computer room has 20 working computers connected through a CAMTEL fibre line, available two afternoons per week; the connection is sometimes slow between 12:00 and 14:00.
\item About 80\% of the students own a smartphone with an MTN or Orange data bundle; bundles are limited, so heavy downloads must be avoided.
\item The school has a Moodle platform hosted on the school server, but it is currently empty.
\item Students already use a free, unsupervised WhatsApp class group of 38 members.
\item Mrs.\ Ngum requires that: (a) every document be traceable (who wrote what, when); (b) no personal data of students (phone numbers, home addresses, photos) be published outside the workspace; (c) the hub must remain usable even if one student leaves the group.
\item The final products are due in \textbf{two weeks}, and each group will present its work in a 10-minute session.
\end{itemize}

\textbf{Your group must deliver:}
\begin{enumerate}
\item A justified choice of collaborative platform;
\item A structured shared folder with one revision document per chapter, roles assigned and permissions correctly set;
\item A co-edited summary sheet of this chapter, with version history enabled;
\item A 10-rule charter of netiquette and security for the whole class;
\item An A3 poster presenting the advantages, the risks encountered, and how the group applied the principles of Law No. 2010/012.
\end{enumerate}

\courssub{Tasks}

\textbf{Task 1 — Analysing the need and choosing the platform.}
\begin{enumerate}
\item List the needs of the class (co-editing, traceability, low data consumption, supervision, permanence of documents) and classify them as \emph{essential} or \emph{desirable}.
\item Compare the three candidate solutions — Google Docs/Drive, the school Moodle space, a supervised WhatsApp or Telegram group — using a comparison table with at least five criteria drawn from the situation's data.
\item Choose one platform (or a justified combination) and defend your choice in five to eight lines, explicitly referring to the constraints (bandwidth, data bundles, supervision, traceability).
\end{enumerate}

\textbf{Task 2 — Building the workspace.}
\begin{enumerate}
\item Design the folder tree of the hub (one folder per subject or per chapter, a folder ``Charter'', a folder ``Past papers'').
\item Assign the four roles — \cle{editor}, \cle{proofreader}, \cle{source-checker}, \cle{moderator} — to the members of your group and state precisely what each role is allowed to do.
\item For each role, indicate the sharing permission you would grant (owner / can edit / can comment / can view) and justify. Explain why the permission ``anyone with the link can edit'' is dangerous here.
\end{enumerate}

\textbf{Task 3 — Co-editing the summary sheet.}
\begin{enumerate}
\item Produce a one-page summary sheet of this chapter (definitions of social network and collaborative environment, types of platforms, privacy settings, netiquette rules, Law No. 2010/012).
\item Describe concretely how your group used real-time co-editing, comments and suggestions, and how the \cle{version history} allowed you to recover an earlier version after a mistaken deletion.
\item Explain how the source-checker verified one piece of information before it was accepted into the sheet.
\end{enumerate}

\textbf{Task 4 — The charter.}
\begin{enumerate}
\item Write a charter of exactly \textbf{10 rules} for the class hub, covering at least: passwords and accounts, respect and netiquette, fake news and verification, protection of personal data, and copyright of shared documents.
\item For two rules of your choice, explain which risk of the situation they prevent.
\end{enumerate}

\textbf{Task 5 — Poster and legal reflection.}
\begin{enumerate}
\item Sketch the layout of your A3 poster: three zones (advantages / risks encountered and solutions / legal framework).
\item State two principles of Law No. 2010/012 (protection of personal data, repression of offences committed through electronic communications, security of information systems) and show concretely how your group applied each of them.
\item Prepare two arguments to answer a sceptical classmate who says: ``Our old WhatsApp group was enough.''
\end{enumerate}

\courssub{Model answer}

\textbf{Task 1 — Choice of platform.}\par
Essential needs: co-editing of documents, traceability of contributions, supervision by the teacher, permanence of files, low data consumption. Desirable: instant messaging, notifications.

\begin{center}
\begin{tabularx}{\linewidth}{|l|X|X|X|}
\hline
\rowcolor{popblueL} \textbf{Criterion} & \textbf{Google Docs/Drive} & \textbf{School Moodle} & \textbf{WhatsApp group} \\
\hline
Real-time co-editing & Excellent & Limited (forums, wikis) & None \\
\hline
Version history / traceability & Built-in, per user & Partial (logs) & None \\
\hline
Supervision by teacher & Easy (teacher owner) & Full (school server) & Weak (38 members, uncontrolled) \\
\hline
Data consumption & Moderate; offline mode possible & Low (text pages) & Low for text, high for media \\
\hline
Permanence of documents & Files stay organised & Files stay organised & Messages buried quickly \\
\hline
Access control & Fine permissions per person & Enrolment-based & Anyone in the group \\
\hline
\end{tabularx}
\end{center}

\textbf{Choice and justification:} We choose \textbf{Google Drive/Docs as the main workspace}, with the teacher as owner, and we keep a \emph{supervised} WhatsApp group only for announcements. Google Docs offers real-time co-editing and a complete version history, which guarantees traceability (requirement (a)); documents remain organised in folders instead of disappearing in a chat flow; and the offline mode plus editing on the school computers two afternoons a week reduces the pressure on students' limited data bundles. Moodle alone would be safer but lacks comfortable real-time co-editing; WhatsApp alone fails traceability and permanence, as the lost summaries and the fake ``leaked paper'' of last term proved.

\textbf{Task 2 — Workspace structure and permissions.}\par
Folder tree: \texttt{RevisionHub\_US6} containing \texttt{ICT/Chapter12}, \texttt{ICT/OtherChapters}, \texttt{Charter}, \texttt{PastPapers}, \texttt{Announcements}.

Roles and permissions:
\begin{itemize}
\item \textbf{Editor} (writes the content): permission \emph{can edit} on the chapter folder;
\item \textbf{Proofreader} (corrects language and layout): \emph{can comment} plus suggestion mode, so corrections remain visible and reversible;
\item \textbf{Source-checker} (verifies facts and references): \emph{can comment}, must approve each fact before the editor integrates it;
\item \textbf{Moderator} (organises, applies the charter, manages members): \emph{can edit} on the whole hub; only Mrs.\ Ngum is \emph{owner}, so the hub survives if a student leaves (requirement (c)).
\end{itemize}
The option ``anyone with the link can edit'' is rejected: a leaked link would let an outsider delete or falsify every summary the night before the mock exam, with no control over membership — exactly the kind of incident the version history can repair but never prevent.

\textbf{Task 3 — Co-editing the summary sheet.}\par
The editor typed the definitions and the typology of platforms; the proofreader worked in \emph{suggestion mode} so that every correction could be accepted or rejected; the source-checker attached a comment to the sentence about Law No. 2010/012 asking for the exact title (``cybersecurity and cybercriminality in Cameroon'') before validation. On the second afternoon, a member accidentally deleted the section on privacy settings: using \emph{File $\rightarrow$ Version history}, the moderator restored the version of 13:42 in less than a minute, and the deleted paragraph was recovered — a practical demonstration that version history is both a traceability tool and a safety net.

\textbf{Task 4 — The 10-rule charter.}
\begin{enumerate}
\item Each member uses a personal account protected by a strong, unique password, never shared, even with friends.
\item Two-factor authentication is activated where available.
\item Every message and document remains respectful: no insults, mockery or personal attacks (netiquette).
\item No information is published before the source-checker has verified it — no fake news, no false ``leaked papers''.
\item No personal data (phone numbers, addresses, photos, results) of any student is shared outside the hub.
\item Only class-related content is posted; chain messages and advertisements are forbidden.
\item Documents copied from a book or website are quoted with their source (copyright).
\item Deletions and major changes are announced in the group before being made.
\item Members log out on shared computers in the school computer room.
\item Any incident (hacked account, suspicious link, conflict) is reported immediately to the moderator and to Mrs.\ Ngum.
\end{enumerate}
Rule 4 directly prevents a repetition of last term's panic caused by the false leaked paper; rule 5 answers requirement (b) of the teacher and the data-protection principle of Law No. 2010/012.

\textbf{Task 5 — Poster and legal reflection.}\par
Poster layout: three vertical zones on A3 paper.

\begin{popfigure}
\begin{tikzpicture}[scale=0.8, every node/.style={font=\small}]
% Zone 1: Advantages
\draw[popblue, thick, rounded corners] (0,0) rectangle (6,10);
\node[popblue, align=center, text width=5.5cm] at (3,9.3) {\textbf{Advantages}};
\node[align=left, text width=5.5cm] at (3,7.5) {
\begin{itemize}
\item Shared workload among 38 students
\item Permanent, organised documents
\item Full traceability (author, date)
\item Accessible from home or computer room
\item Offline mode for low bandwidth
\end{itemize}
};
% Zone 2: Risks & Solutions
\draw[poporange, thick, rounded corners] (6.5,0) rectangle (12.5,10);
\node[poporange, align=center, text width=5.5cm] at (9.5,9.3) {\textbf{Risks \& Solutions}};
\node[align=left, text width=5.5cm] at (9.5,7.5) {
\begin{itemize}
\item Accidental deletion $\rightarrow$ Version history
\item Slow connection at midday $\rightarrow$ Offline work + lab sessions
\item Fake news / leaked papers $\rightarrow$ Source-checker role
\item Account theft $\rightarrow$ Rules 1, 2, 10 (2FA, reporting)
\item Data leakage $\rightarrow$ Restricted permissions, no personal data
\end{itemize}
};
% Zone 3: Legal Framework
\draw[popgreen, thick, rounded corners] (13,0) rectangle (19,10);
\node[popgreen, align=center, text width=5.5cm] at (16,9.3) {\textbf{Law No. 2010/012}};
\node[align=left, text width=5.5cm] at (16,7.5) {
\begin{itemize}
\item \textbf{Art. 7-8: Personal data protection} $\rightarrow$ No student data published; teacher owns workspace; access restricted to named members.
\item \textbf{Art. 12-15: Security of information systems} $\rightarrow$ Strong passwords, 2FA, immediate incident reporting (Rules 1,2,10).
\item \textbf{Art. 20-25: Cyber offences} $\rightarrow$ Charter forbids fake news, identity theft, unauthorised access; sanctions apply.
\end{itemize}
};
\end{tikzpicture}
\legende{Proposed A3 poster layout: three coloured zones for Advantages, Risks \& Solutions, and Legal Framework.}
\end{popfigure}

Two principles of Law No. 2010/012 applied concretely:
\begin{itemize}
\item \textbf{Protection of personal data and privacy in electronic communications (Articles 7-8):} we restricted access to named members, banned the publication of students' personal data (rule 5) and kept the teacher as owner of the space.
\item \textbf{Security of information systems and repression of cyber-offences (Articles 12-15, 20-25):} strong passwords, two-factor authentication, immediate reporting of incidents (rules 1, 2, 10) protect the workspace against intrusion, identity theft and the deliberate diffusion of false information, which the law punishes.
\end{itemize}

Answer to the sceptical classmate: first, WhatsApp keeps no organised, recoverable record — last term's summaries were lost and the fake paper spread unchecked, whereas our hub keeps every version with its author and date; second, our hub separates discussion from documents and is supervised, so it satisfies the teacher's traceability and data-protection requirements, which an uncontrolled 38-member group never could.

\begin{attention}[Assessment criteria]
Mrs.\ Ngum will evaluate: the relevance of the platform choice, the correctness of permissions and roles, the effective use of version history, the completeness and realism of the charter, the accuracy of the legal references, and the quality of the oral presentation. A beautiful poster with wrong permissions scores less than a simple but correctly configured workspace.
\end{attention}

\newpage

\coursec{Methods and worked examples}

\courssub{Methods and Worked Examples}

\begin{methode}[Creating and Securing a Social Network Account]
\begin{enumerate}
    \item \textbf{Choose the right platform} for your purpose: professional (LinkedIn), general (Facebook, WhatsApp), microblogging (X/Twitter), media sharing (Instagram, TikTok, YouTube). In Cameroon, WhatsApp and Facebook dominate daily communication; LinkedIn is standard for job hunting.
    \item \textbf{Register with a strong, unique password}: at least 12 characters, mixing upper/lower case, digits, symbols. Avoid personal data (name, birth date, phone number). Use a passphrase: \texttt{Mtn@Yaounde2025!Secure}.
    \item \textbf{Enable Two-Factor Authentication (2FA)} immediately: prefer an authenticator app (Google Authenticator, Microsoft Authenticator) over SMS, because SIM-swap fraud is common in Cameroon.
    \item \textbf{Set up recovery options}: add a secondary email and a trusted phone number (your active MTN/Orange line). Verify them.
    \item \textbf{Review default privacy settings} before posting: who can see future posts, who can tag you, profile visibility to search engines.
    \item \textbf{Log out of unused devices} regularly: Settings $\rightarrow$ Security $\rightarrow$ Where you're logged in $\rightarrow$ Log out of sessions you don't recognise.
    \item \textbf{Keep the app and OS updated} to patch vulnerabilities.
\end{enumerate}
\end{methode}

\begin{exoresolu}[Securing a New Facebook Account for a Small Business in Douala]
    \textbf{Statement:} Marie has just opened a Facebook page for her \emph{ndolé} delivery service in Douala. She uses her personal phone number (MTN) and email \texttt{marie.ndole@yahoo.fr}. She wants to protect the page from hijacking and limit who can post on the timeline. Describe, step by step, the security and privacy settings she should configure immediately after creating the account.
    \tcblower
    \textbf{Detailed Solution:}
    \begin{enumerate}
        \item \textbf{Strong password:} Marie goes to \textit{Settings \& Privacy $\rightarrow$ Settings $\rightarrow$ Security and Login $\rightarrow$ Change password}. She creates a 14-character passphrase: \texttt{Douala-Ndole\#2025!Mtn}. She does \emph{not} reuse her Yahoo or Mobile Money PIN.
        \item \textbf{Two-Factor Authentication (2FA):} Same menu $\rightarrow$ \textit{Use two-factor authentication} $\rightarrow$ Choose \textit{Authentication app} (Google Authenticator). She scans the QR code, enters the 6-digit code. \textbf{Why not SMS?} SIM-swap attacks are frequent in Cameroon; an app generates codes offline.
        \item \textbf{Trusted contacts:} \textit{Setting up extra security $\rightarrow$ Choose 3 to 5 friends} who can help recover the account if locked out. She picks reliable family members.
        \item \textbf{Login alerts:} \textit{Get alerts about unrecognised logins} $\rightarrow$ Enable both \textit{Notifications} and \textit{Email} alerts. She will know instantly if someone logs in from Yaoundé or abroad.
        \item \textbf{Page roles (critical for business):} On the Page $\rightarrow$ \textit{Settings $\rightarrow$ Page Roles}. She assigns herself as \textit{Admin}, her sister as \textit{Editor} (can post but not change roles), and \textbf{never} gives Admin to an external marketer.
        \item \textbf{Visitor posts:} \textit{Settings $\rightarrow$ General $\rightarrow$ Visitor Posts} $\rightarrow$ \textit{Disable posts by other people on the Page} or \textit{Review posts by other people before they are published}. This prevents spam or malicious links on her business timeline.
        \item \textbf{Profile privacy (personal profile linked to page):} \textit{Settings $\rightarrow$ Privacy $\rightarrow$ Your Activity} $\rightarrow$ \textit{Who can see your future posts?} $\rightarrow$ \textit{Friends} (not Public). \textit{Limit past posts} to restrict old public posts.
        \item \textbf{Search engine visibility:} \textit{Settings $\rightarrow$ Privacy $\rightarrow$ How people find and contact you} $\rightarrow$ \textit{Do you want search engines outside Facebook to link to your profile?} $\rightarrow$ \textit{No} (for personal profile; the Page itself should remain public for customers).
        \item \textbf{Device management:} \textit{Security and Login $\rightarrow$ Where you're logged in} $\rightarrow$ Log out of any old phones or cybercafé computers.
        \item \textbf{Updates:} She enables auto-update for the Facebook app via Google Play Store / App Store and keeps Android/iOS updated.
    \end{enumerate}
    \textbf{Key reflex:} Treat the business page like a bank account — only trusted people have Admin rights, 2FA is mandatory, and login alerts are your early warning system.
\end{exoresolu}

\begin{methode}[Managing Privacy Settings and Digital Footprint]
\begin{enumerate}
    \item \textbf{Audit your profile visibility}: Use the platform's \textit{View As} tool (Facebook) or \textit{Privacy Checkup} to see what strangers see.
    \item \textbf{Apply the principle of least privilege}: Set default audience for new posts to \textit{Friends} or a custom list (e.g., \textit{Close Friends}, \textit{Work}), never \textit{Public} by default.
    \item \textbf{Control tagging}: Enable \textit{Review tags people add to your posts} and \textit{Review posts you're tagged in before they appear on your timeline}.
    \item \textbf{Limit past posts}: Use the \textit{Limit the audience for posts you've shared with friends of friends or Public?} tool to retroactively restrict old posts.
    \item \textbf{Manage third-party apps}: Remove apps/websites you no longer use that have access to your profile (Settings $\rightarrow$ Apps and Websites).
    \item \textbf{Clean your digital footprint periodically}: Search your name on Google, DuckDuckGo, and Cameroonian directories; request removal of outdated or sensitive info (old phone numbers, school reports, etc.).
    \item \textbf{Separate personal and professional identities}: Use a dedicated professional profile/page for work; keep personal life on a restricted personal profile.
\end{enumerate}
\end{methode}

\begin{exoresolu}[Auditing and Reducing Digital Footprint for a Job Seeker in Yaoundé]
    \textbf{Statement:} Paul, a final-year student at the University of Yaoundé I, is applying for internships at SONARA and Orange Cameroon. He realises his Facebook profile (public since 2018) contains party photos, political comments, and his personal phone number. He also used his Facebook login for several quiz apps. Describe a complete cleanup procedure to present a professional digital footprint.
    \tcblower
    \textbf{Detailed Solution:}
    \begin{enumerate}
        \item \textbf{Privacy Checkup (Facebook):} \textit{Settings \& Privacy $\rightarrow$ Privacy Checkup}. He goes through the four steps:
              \begin{itemize}
                  \item \textit{Who can see what you share} $\rightarrow$ Future posts: \textit{Friends}; Profile info (phone, email, birthday): \textit{Only me}; Work/Education: \textit{Friends}.
                  \item \textit{How to keep your account secure} $\rightarrow$ Already done 2FA (see previous method).
                  \item \textit{How people can find you} $\rightarrow$ Friend requests: \textit{Friends of friends}; Phone/Email lookup: \textit{Friends}; Search engines: \textit{Off}.
                  \item \textit{Your data settings on Facebook} $\rightarrow$ Apps and Websites: \textit{Remove all quiz apps, games, and unused logins}.
              \end{itemize}
        \item \textbf{Limit Past Posts:} \textit{Settings $\rightarrow$ Privacy $\rightarrow$ Your Activity $\rightarrow$ Limit Past Posts} $\rightarrow$ \textit{Limit Past Posts} button. This changes all previous \textit{Public} and \textit{Friends of Friends} posts to \textit{Friends} only. \textbf{Warning:} This is irreversible in bulk; he cannot easily make them public again one by one.
        \item \textbf{Timeline and Tagging:} \textit{Settings $\rightarrow$ Privacy $\rightarrow$ Timeline and Tagging} $\rightarrow$ Enable \textit{Review posts you're tagged in before they appear on your timeline?} $\rightarrow$ \textit{Yes}. Enable \textit{Review tags people add to your posts?} $\rightarrow$ \textit{Yes}.
        \item \textbf{Profile Cleanup (manual):} He uses \textit{View As} (three dots on profile $\rightarrow$ \textit{View As}) to see what a stranger sees. He manually hides/deletes specific embarrassing photos/posts (three dots on post $\rightarrow$ \textit{Move to Archive} or \textit{Delete}).
        \item \textbf{Google Search Audit:} He searches \texttt{"Paul Atanga" Yaoundé}, \texttt{"Paul Atanga" Facebook}, \texttt{"Paul Atanga" phone}. He notes any unwanted results (old blog comments, public event pages). For each, he either deletes the source or uses Google's \textit{Remove outdated content} tool / \textit{Right to be forgotten} request (though Cameroon's law is evolving, Google honours EU-style requests globally).
        \item \textbf{Create Professional Presence:} He creates a \textbf{LinkedIn} profile with a professional photo, headline: \textit{Final-year Engineering Student | Petroleum \& Gas | Seeking Internship}, detailed education, skills, and a summary. He sets LinkedIn visibility to \textit{Public} for recruiters.
        \item \textbf{Separate Identities:} He converts his old Facebook profile to a \textit{Page} (if he wants to keep followers) or keeps it strictly private (\textit{Friends only}) and uses the new LinkedIn for all professional networking. He adds the LinkedIn URL to his CV.
        \item \textbf{Ongoing Habit:} Before any post, he asks: \textit{``Would I want my future boss at Orange to see this?''} If no $\rightarrow$ don't post, or post to a \textit{Close Friends} list only.
    \end{enumerate}
    \textbf{Examiner's tip:} The \textit{Limit Past Posts} feature is a frequent GCE question — know it changes audience retroactively but not the posts themselves.
\end{exoresolu}

\begin{methode}[Identifying Types of Social Networks and Their Use Cases]
\begin{enumerate}
    \item \textbf{Classify by primary function}:
          \begin{itemize}
              \item \textit{Social networking (general)}: Facebook, WhatsApp — connecting people, sharing life updates.
              \item \textit{Professional networking}: LinkedIn, Viadeo — careers, B2B, recruitment.
              \item \textit{Microblogging / Real-time}: X (Twitter), Threads — news, opinions, short updates.
              \item \textit{Media sharing}: Instagram, TikTok, YouTube, Snapchat — photos, short/long videos.
              \item \textit{Discussion forums / Communities}: Reddit, Quora, Nairaland (Nigeria-focused but used in Cameroon) — Q\&A, niche interests.
              \item \textit{Instant messaging / Collaboration}: WhatsApp, Telegram, Signal, Microsoft Teams, Slack — private/group chat, file sharing, calls.
          \end{itemize}
    \item \textbf{Match platform to objective}:
          \begin{itemize}
              \item \textit{Job hunting in Cameroon} $\rightarrow$ LinkedIn (formal CV), Facebook groups (e.g., \textit{Emploi Cameroun}, \textit{Job Alert Cameroon}), WhatsApp broadcast lists from recruitment agencies.
              \item \textit{Marketing a local fashion brand} $\rightarrow$ Instagram (visuals), TikTok (short videos/reels), WhatsApp Business (catalogue, orders, Mobile Money payment links).
              \item \textit{Citizen journalism / Breaking news} $\rightarrow$ X (hashtags like \#Cameroun, \#Yaounde), Facebook Live.
              \item \textit{Team project coordination} $\rightarrow$ Microsoft Teams / Google Workspace (documents, meetings), WhatsApp group (quick alerts).
          \end{itemize}
    \item \textbf{Consider audience demographics in Cameroon}: Facebook/WhatsApp cover all ages; TikTok/Instagram skew younger (15–30); LinkedIn is urban, educated, formal sector; X is journalists, activists, tech.
    \item \textbf{Evaluate features}: Groups/Pages/Events (Facebook), Channels/Groups (Telegram), Communities (WhatsApp), Spaces (X), Hashtags/Trends (all).
\end{enumerate}
\end{methode}

\begin{exoresolu}[Choosing the Right Platforms for a Youth NGO in Buea]
    \textbf{Statement:} \textit{Green Youth Cameroon}, an NGO based in Buea, wants to: (a) recruit volunteers aged 18–25, (b) showcase tree-planting events with photos/videos, (c) coordinate logistics among 15 core members, (d) advocate for policy change to the Ministry of Environment (MINEPDED). Recommend the most suitable platform(s) for each objective, justifying with features and Cameroonian context.
    \tcblower
    \textbf{Detailed Solution:}
    \begin{enumerate}
        \item \textbf{Recruit volunteers (18–25)}:
              \begin{itemize}
                  \item \textbf{Primary: TikTok \& Instagram Reels.} Short, engaging videos of past events (planting, nursery work) with trending sounds reach the youth demographic algorithmically. Hashtags: \#Buea, \#CameroonYouth, \#ClimateAction, \#GreenCameroon.
                  \item \textbf{Secondary: Facebook Groups.} Create a public group \textit{Green Youth Cameroon Volunteers} for announcements, Q\&A, and peer interaction. Many Cameroonian students still use Facebook heavily for groups.
                  \item \textbf{Why not LinkedIn?} Too formal; 18–25 age group in Cameroon rarely maintains active LinkedIn profiles.
              \end{itemize}
        \item \textbf{Showcase events (photos/videos)}:
              \begin{itemize}
                  \item \textbf{Instagram (Feed + Highlights + Reels)}: Curated grid for best photos; Highlights for each event (e.g., \textit{Mount Cameroon Planting 2024}); Reels for viral reach.
                  \item \textbf{YouTube Channel}: Longer documentaries (5–10 min) for donors/partners, embedded on website.
                  \item \textbf{Facebook Page}: Album feature for high-res photo downloads by press/partners; Events feature for past/future events.
              \end{itemize}
        \item \textbf{Coordinate 15 core members}:
              \begin{itemize}
                  \item \textbf{WhatsApp Group (with Communities)}: Universal in Cameroon; low data usage; voice notes for quick updates; file sharing (PDF budgets, GPS locations). Use \textit{Communities} to separate \textit{Logistics}, \textit{Finance}, \textit{Communications} sub-groups under one umbrella.
                  \item \textbf{Google Workspace (free for non-profits)}: Shared Drive for documents (proposals, reports), Google Meet for monthly virtual meetings, Calendar for scheduling. Apply for \textit{Google for Nonprofits} (requires registration with MINAS/Ministry of Youth).
              \end{itemize}
        \item \textbf{Advocate to MINEPDED}:
              \begin{itemize}
                  \item \textbf{X (Twitter)}: Tag \texttt{@MINEPDED\_CM}, \texttt{@Minjeunes\_CM}, \texttt{@PRC\_CM}. Use hashtags \#CamerounVert, \#COP28. Journalists and ministry communication officers monitor X.
                  \item \textbf{LinkedIn}: Publish a formal \textit{Policy Brief} article tagging the Minister's LinkedIn profile (if exists) and relevant officials. Professional tone, PDF attachment.
                  \item \textbf{Email + Physical Courtesy Visit}: Digital complements, not replaces, formal administrative channels in Cameroon.
              \end{itemize}
    \end{enumerate}
    \textbf{Summary Table:}
    \begin{center}
    \begin{tabularx}{\linewidth}{|l|X|X|}\hline
    \rowcolor{popblueL} \textbf{Objective} & \textbf{Primary Platform(s)} & \textbf{Key Feature Used} \\ \hline
    Recruit volunteers & TikTok, Instagram, Facebook Group & Algorithm reach, Groups \\ \hline
    Showcase events & Instagram, YouTube, Facebook Page & Reels/Highlights, Albums, Events \\ \hline
    Coordinate team & WhatsApp Community, Google Workspace & Sub-groups, Shared Drive, Meet \\ \hline
    Policy advocacy & X, LinkedIn, Formal Email & Tagging officials, Articles \\ \hline
    \end{tabularx}
    \end{center}
\end{exoresolu}

\begin{methode}[Using Collaborative Environments for Teamwork (Google Workspace / Microsoft 365)]
\begin{enumerate}
    \item \textbf{Choose the ecosystem}: Google Workspace (Drive, Docs, Sheets, Slides, Meet, Calendar) — cloud-native, real-time co-editing, free for personal/education. Microsoft 365 (OneDrive, Word, Excel, PowerPoint, Teams) — desktop-app integration, enterprise standard in Cameroon (many ministries, banks, SONARA use it).
    \item \textbf{Set up a shared structure}:
          \begin{itemize}
              \item Create a \textbf{Shared Drive (Google) / Team Site (Teams/SharePoint)} owned by the organisation, not an individual.
              \item Folder hierarchy: \texttt{/Project\_Name/01\_Admin/02\_Research/03\_Drafts/04\_Final/05\_Assets}.
              \item Set permissions: \textit{Content Manager} (edit/move/delete) for core team; \textit{Commenter} for reviewers; \textit{Viewer} for stakeholders.
          \end{itemize}
    \item \textbf{Work on documents simultaneously}: Open the file in the browser (Docs/Sheets/Slides or Word/Excel/PowerPoint Online). See coloured cursors. Use \textit{Comments} (@mention teammates) and \textit{Suggesting/Track Changes} mode for reviews.
    \item \textbf{Version history}: \textit{File $\rightarrow$ Version history $\rightarrow$ Name current version} (e.g., \textit{v1.0\_Submission\_MINEPDED}). Restore if needed.
    \item \textbf{Communicate in context}: Use \textit{Teams Channels} or \textit{Google Chat Spaces} linked to the project — keeps discussion attached to files, not lost in WhatsApp.
    \item \textbf{Schedule and run meetings}: Calendar $\rightarrow$ Create event $\rightarrow$ Add video conferencing (Meet/Teams) $\rightarrow$ Attach agenda doc. Record meeting (with consent) $\rightarrow$ auto-saves to Drive/OneDrive.
    \item \textbf{Offline access}: Enable \textit{Offline mode} in Drive/OneDrive settings for critical files when travelling (e.g., fieldwork in Far North with poor 4G).
\end{enumerate}
\end{methode}

\begin{exoresolu}[Setting Up a Collaborative Workspace for a Final-Year Project Group]
    \textbf{Statement:} A group of four Upper Sixth students (Awa, Boris, Cindy, Didier) from Government Bilingual High School Molyko must write a 30-page research project on \textit{``Mobile Money Adoption in Rural Southwest Region''}. They have one month. They have mixed devices (Android, iPhone, Windows laptop, no laptop).Design a complete collaborative workflow using free tools accessible in Cameroon.
    \tcblower
    \textbf{Detailed Solution:}
    \begin{enumerate}
        \item \textbf{Platform Choice: Google Workspace (free, works on all devices via browser, low bandwidth for editing).}
        \item \textbf{Account Setup:} Each creates a Google account (or uses existing). One (Awa) creates a \textbf{Shared Drive} named \textit{GBHS\_Molyko\_Project\_MobileMoney}. She adds the other three as \textit{Content Managers}.
        \item \textbf{Folder Structure (created by Awa):}
              \begin{verbatim}
    /GBHS\_Molyko\_Project\_MobileMoney
      ├── 01\_Admin
      │     ├── Timeline\_Gantt.xlsx
      │     ├── Meeting\_Minutes.docx
      │     └── Task\_Assignment.xlsx
      ├── 02\_Literature\_Review
      │     ├── Papers\_To\_Read/   (PDFs)
      │     └── Synthesis.docx
      ├── 03\_Data\_Collection
      │     ├── Questionnaire\_Google\_Form/  (link)
      │     ├── Raw\_Data\_CSV/
      │     └── Field\_Notes\_Photos/
      ├── 04\_Analysis
      │     ├── SPSS\_Output/   (or Sheets)
      │     └── Charts\_Graphs/
      ├── 05\_Drafting
      │     ├── Chapter$\_{1}$\_Intro.docx
      │     ├── Chapter$\_{2}$\_LitRev.docx
      │     ├── Chapter$\_{3}$\_Methodology.docx
      │     ├── Chapter$\_{4}$\_Results.docx
      │     ├── Chapter$\_{5}$\_Discussion.docx
      │     └── Full\_Draft\_Master.docx
      ├── 06\_Final\_Submission
      │     ├── Final\_PDF.pdf
      │     ├── Plagiarism\_Report.pdf
      │     └── Defense\_Slides.pptx
      └── 07\_Assets
            ├── Logo\_School.png
            └── References\_BibTeX.bib
              \end{verbatim}
        \item \textbf{Workflow Rules (agreed in first WhatsApp call):}
              \begin{itemize}
                  \item \textbf{Naming convention:} \texttt{YYYY-MM-DD\_Initials\_Description\_vN.ext} (e.g., \texttt{2025-01-15\_BC\_LitRev\_Summary\_v2.docx}).
                  \item \textbf{Editing:} Only edit files \textbf{inside} the Shared Drive (browser). Never download, edit offline, re-upload (creates conflicts).
                  \item \textbf{Comments:} Use \textit{Comment} (@name) for questions; \textit{Suggesting mode} for edits. Resolve comments when done.
                  \item \textbf{Weekly Sync:} Sunday 20:00 on Google Meet (link in Calendar event attached to \textit{01\_Admin/Meeting\_Minutes.docx}). Record for absent members.
                  \item \textbf{Backup:} Awa enables \textit{Drive for Desktop} on her laptop for local sync; others rely on cloud. Monthly, Awa downloads the whole Drive as ZIP to her laptop/USB.
              \end{itemize}
        \item \textbf{Key Tools Used:}
              \begin{itemize}
                  \item \textbf{Google Forms:} Questionnaire (offline-capable via app) $\rightarrow$ responses auto-populate \textit{03\_Data\_Collection/Raw\_Data\_CSV}.
                  \item \textbf{Google Sheets:} Data cleaning, pivot tables, charts $\rightarrow$ \textit{04\_Analysis}.
                  \item \textbf{Google Docs:} All writing. \textit{Tools $\rightarrow$ Citations} for APA/MLA. \textit{Add-ons $\rightarrow$ DocuSign} for supervisor signature (if needed).
                  \item \textbf{Google Slides:} Defense presentation — real-time co-design.
                  \item \textbf{Google Calendar:} Shared calendar \textit{Project Deadlines} with reminders (1 week, 1 day before each milestone).
              \end{itemize}
        \item \textbf{Handling Low Connectivity (Didier, no laptop, poor 4G in village):}
              \begin{itemize}
                  \item He installs \textit{Google Docs, Sheets, Slides apps} on his phone $\rightarrow$ enables \textit{Make recent files available offline} in each app.
                  \item He downloads the current chapter he's working on before travelling.
                  \item He writes/edits offline; changes sync automatically when he returns to 4G/WiFi zone (e.g., school computer lab).
                  \item Voice typing (mic icon in Docs mobile) helps write long sections without typing on phone keyboard.
              \end{itemize}
        \item \textbf{Finalisation:} \textit{File $\rightarrow$ Version history $\rightarrow$ Name version ``v1.0\_Submission\_2025-02-20''}. Export \textit{Final\_PDF.pdf} (File $\rightarrow$ Download $\rightarrow$ PDF). Run plagiarism check (free: \textit{Quetext}, \textit{SmallSEOTools}) $\rightarrow$ save report in \textit{06\_Final\_Submission}.
    \end{enumerate}
    \textbf{Examiner's note:} GCE questions often ask for \textit{advantages of cloud collaboration over email attachments} — answer: real-time co-editing, single source of truth, version history, access from any device, no attachment size limits.
\end{exoresolu}

\begin{methode}[Applying Netiquette and Recognising Cyberbullying / Online Harassment]
\begin{enumerate}
    \item \textbf{Core Netiquette Rules (adapted for Cameroon context)}:
          \begin{enumerate}
              \item \textbf{Respect the human}: Behind every screen is a person. No insults, tribal slurs (e.g., \textit{Anglophone/Francophone} hate), religious mockery, body shaming.
              \item \textbf{Write clearly, avoid ALL CAPS} (perceived as shouting). Use proper French/English/Pidgin; limit abbreviations in formal contexts.
              \item \textbf{Context matters}: WhatsApp family group $\neq$ LinkedIn comment $\neq$ X political thread. Adapt tone.
              \item \textbf{Verify before sharing}: \textit{``Forwarded many times''} label on WhatsApp = red flag. Check with \textit{Africa Check}, \textit{DataCheck}, official ministry websites before forwarding health/political/security alerts.
              \item \textbf{Respect privacy}: Don't share screenshots of private chats, photos, or documents without consent. Don't add people to groups without asking.
              \item \textbf{Give credit}: Share original source (link, handle) for images, videos, articles. Don't crop watermarks.
              \item \textbf{Disagree constructively}: Attack the argument, not the person. \textit{``I disagree because data from INS shows...''} not \textit{``You're stupid/biased.''}
          \end{enumerate}
    \item \textbf{Identify Cyberbullying / Harassment}:
          \begin{itemize}
              \item Repeated, intentional, power imbalance.
              \item Forms: \textit{Harassment} (repeated insults/threats), \textit{Doxing} (publishing private info: address, phone, ID), \textit{Revenge porn} (non-consensual intimate images), \textit{Impersonation} (fake profile), \textit{Exclusion} (deliberate blocking from groups), \textit{Cyberstalking}.
          \end{itemize}
    \item \textbf{Response Protocol (STOP)}:
          \begin{enumerate}
              \item \textbf{S}creenshot evidence (timestamp, URL, profile).
              \item \textbf{T}ell a trusted adult / school authority / HR (if workplace).
              \item \textbf{O}pt out: Block, restrict, report on platform.
              \item \textbf{P}reserve mental health: Talk to counsellor; don't engage with trolls.
          \end{enumerate}
    \item \textbf{Reporting Channels in Cameroon}:
          \begin{itemize}
              \item Platform reporting tools (Facebook/WhatsApp/Instagram/X reporting flows).
              \item \textbf{ANTIC (Agence Nationale des Technologies de l'Information et de la Communication)}: Cybercrime reporting portal \url{https://antic.cm} or toll-free \texttt{8202} (MTN/Orange). They handle cyberbullying, sextortion, fake news, identity theft.
              \item \textbf{Police / Gendarmerie}: For threats of physical violence, sextortion, fraud — file a \textit{procès-verbal}.
              \item \textbf{MINAS (Ministry of Social Affairs)}: Child protection unit for minors.
          \end{itemize}
\end{enumerate}
\end{methode}

\begin{exoresolu}[Handling a Cyberbullying Incident in a School WhatsApp Group]
    \textbf{Statement:} In a WhatsApp group \textit{``Upper Sixth Science 2025 - GBHS Bamenda''} (120 students), a student repeatedly posts edited photos mocking another student's physical disability, with captions in Pidgin. The victim has stopped attending school. As the class prefect, describe the immediate actions you take, the netiquette rules violated, and the reporting channels available in Cameroon.
    \tcblower
    \textbf{Detailed Solution:}
    \begin{enumerate}
        \item \textbf{Immediate Actions (STOP Protocol):}
              \begin{enumerate}
                  \item \textbf{Screenshot:} I take screenshots of \textbf{every} offensive message (photo + caption + sender's name/number + timestamp + group name). I save them in a secure folder on my phone and email them to myself (backup).
                  \item \textbf{Support Victim:} I privately message the victim: \textit{``I saw what happened. It's not your fault. I'm taking action. Do you want me to involve the discipline master / your parents?''} I do \emph{not} confront the bully in the group.
                  \item \textbf{Group Management:} As admin (or asking the admin), I \textbf{remove the bully} from the group immediately. I post a \textbf{single, calm message}: \textit{``This group is for academic support. Harassment and mockery of any student are prohibited. The offending member has been removed. Further incidents will be reported to school administration and ANTIC.''} I then \textbf{restrict group settings}: \textit{Group Settings $\rightarrow$ Send Messages $\rightarrow$ Only Admins} (temporarily) to stop pile-on.
                  \item \textbf{Report on WhatsApp:} Long-press each offensive message $\rightarrow$ \textit{Report} $\rightarrow$ \textit{Report and Block}. This sends the last 5 messages to WhatsApp moderation.
              \end{enumerate}
        \item \textbf{Netiquette Rules Violated:}
              \begin{itemize}
                  \item \textbf{Rule 1: Respect the human} — mocking a disability is dehumanising.
                  \item \textbf{Rule 4: Verify before sharing} — the bully created/shares doctored images (disinformation + harassment).
                  \item \textbf{Rule 5: Respect privacy} — using the victim's image without consent.
                  \item \textbf{Rule 7: Disagree constructively} — this isn't disagreement; it's targeted abuse.
                  \item \textbf{Context:} A school WhatsApp group is a \textit{semi-formal educational space}; such behaviour violates school regulations and Cameroonian law.
              \end{itemize}
        \item \textbf{Reporting Channels in Cameroon:}
              \begin{enumerate}
                  \item \textbf{School Administration:} Discipline Master $\rightarrow$ Principal. They can sanction under school rules (suspension, expulsion) and summon parents.
                  \item \textbf{ANTIC:} I file a report on \url{https://antic.cm} (or call \texttt{8202}) with: \textit{Category: Cyberbullying / Harassment}; \textit{Evidence: screenshots, group link, phone numbers}; \textit{Victim details (with consent)}. ANTIC can request WhatsApp/Meta to take down content and identify the perpetrator (via IP/subscriber info from MTN/Orange).
                  \item \textbf{Gendarmerie / Police (Brigade de Recherche)}: If the bullying includes threats of violence or sextortion, the victim's parents should file a \textit{plainte} (formal complaint). Cybercrime law (Law No. 2010/012) penalises \textit{``offences committed by means of electronic communication''} (Art. 77-79).
                  \item \textbf{MINAS / MINPROFF:} If the victim is a minor ($<18$), the Ministry of Social Affairs / Ministry of Women's Empowerment and the Family can intervene for psychosocial support.
              \end{enumerate}
        \item \textbf{Follow-up:}
              \begin{itemize}
                  \item Monitor group for retaliation / secondary victimisation.
                  \item Organise a class talk (with counsellor) on cyberbullying, empathy, and digital citizenship.
                  \item Update group description with \textit{``Zero tolerance for bullying. Report to admin/ANTIC 8202.''}
              \end{itemize}
    \end{enumerate}
    \textbf{Key distinction:} \textit{Netiquette} = social norms (soft skills); \textit{Cybercrime Law} = legal sanctions (hard consequences). Both apply.
\end{exoresolu}

\begin{methode}[Understanding the Legal Framework: Cameroon Cybersecurity and Data Protection]
\begin{enumerate}
    \item \textbf{Key Laws (Cameroon)}:
          \begin{itemize}
              \item \textbf{Law No. 2010/012 of 21 December 2010} on Cybersecurity and Cybercrime. Criminalises: illegal access to systems, data interception, system interference, computer-related forgery/fraud, child pornography, \textbf{offensive communication} (Art. 78: \textit{``whoever uses electronic communications to insult, threaten, harass...''}), identity theft.
              \item \textbf{Law No. 2024/014 of 24 December 2024} on the Protection of Personal Data (new, aligns with AU Convention / GDPR principles). Establishes \textbf{ANTIC} as Data Protection Authority. Principles: consent, purpose limitation, data minimisation, accuracy, storage limitation, integrity/confidentiality, accountability. Rights: access, rectification, erasure (right to be forgotten), objection, portability.
              \item \textbf{Law No. 2010/013 of 21 December 2010} on Electronic Communications (regulates operators, licensing, universal service).
              \item \textbf{Decree No. 2013/0399/PM} on the organisation of ANTIC.
          \end{itemize}
    \item \textbf{Key Concepts for GCE}:
          \begin{itemize}
              \item \textbf{Personal Data}: Any info relating to an identified/identifiable natural person (name, ID, phone, IP, biometric, location).
              \item \textbf{Sensitive Data}: Racial/ethnic origin, political opinions, religious beliefs, health, genetic, biometric, sexual orientation — requires \textbf{explicit consent}.
              \item \textbf{Data Controller} vs \textbf{Data Processor}: Controller decides \textit{why} and \textit{how} (e.g., a bank); Processor acts on controller's behalf (e.g., cloud provider).
              \item \textbf{Data Subject Rights}: You can ask MTN/Orange/Facebook/your school: \textit{``What data do you hold on me? Delete it. Give me a copy.''}
              \item \textbf{Breach Notification}: Controller must notify ANTIC and affected subjects within \textbf{72 hours} of discovering a breach (Art. 33 GDPR-style, now in Cameroon law).
          \end{itemize}
    \item \textbf{Practical Compliance for Students/Users}:
          \begin{itemize}
              \item Don't share others' personal data (photos, phone numbers, exam results) without consent.
              \item Use privacy settings (previous methods) to exercise your rights.
              \item Report data breaches (e.g., your school leaks student list) to ANTIC.
              \item Know that \textit{``freedom of expression''} does not protect hate speech, fake news, or harassment (Law 2010/012 Art. 78).
          \end{itemize}
\end{enumerate}
\end{methode}

\begin{exoresolu}[Applying Data Protection Law to a School's Digital Practices]
    \textbf{Statement:} A private high school in Yaoundé uses a WhatsApp group per class to share: (a) students' full names, admission numbers, and termly grades (PDF), (b) photos of students during sports day (faces visible), (c) parents' phone numbers (visible to all group members). The school also stores all student records on a Google Drive folder shared with \textit{``Anyone with the link can view''}. Analyse the compliance with Cameroon's Law No. 2024/014 on Personal Data Protection and propose corrective measures.
    \tcblower
    \textbf{Detailed Solution:}
    \begin{enumerate}
        \item \textbf{Legal Analysis (Law No. 2024/014):}
              \begin{itemize}
                  \item \textbf{Data Controller:} The school (represented by the Principal) determines purposes/means of processing student/parent data.
                  \item \textbf{Personal Data Processed:}
                        \begin{itemize}
                            \item \textit{Grades + Admission numbers} $\rightarrow$ \textbf{Personal data} (academic records).
                            \item \textit{Photos (faces)} $\rightarrow$ \textbf{Biometric data} (special category / sensitive data under Art. 9 GDPR-style, likely sensitive in Cameroon law too).
                            \item \textit{Parents' phone numbers} $\rightarrow$ \textbf{Personal data} (contact details).
                        \end{itemize}
                  \item \textbf{Lawful Basis?} \textit{Contract (education service)} or \textit{Legal obligation (MINEDUB reporting)} may cover grades. \textbf{But:} \textit{Consent} is required for photos (sensitive data) and for sharing phone numbers with third parties (other parents).
                  \item \textbf{Violations Identified:}
                        \begin{enumerate}
                            \item \textbf{No Consent for Photos:} Posting children's faces on WhatsApp (accessible to all parents, forwardable) without explicit, informed, written consent from parents/guardians. Violation of \textit{purpose limitation} and \textit{sensitive data} rules.
                            \item \textbf{Excessive Disclosure (Grades):} Sharing \textit{all} students' grades with \textit{all} parents violates \textit{data minimisation} and \textit{confidentiality}. Each parent should only access their own child's data.
                            \item \textbf{Phone Number Exposure:} Sharing parents' numbers with 50+ strangers (other parents) without consent violates \textit{privacy by design} and \textit{security of processing}.
                            \item \textbf{Insecure Storage (Google Drive):} \textit{``Anyone with the link can view''} = \textbf{public access}. No access control, no encryption at rest (beyond Google's), no audit trail. Violates \textit{integrity and confidentiality} (Art. 32 GDPR / equivalent). High risk of data breach.
                            \item \textbf{No Data Protection Officer (DPO):} School likely hasn't appointed a DPO (mandatory for systematic monitoring of minors).
                            \item \textbf{No Privacy Notice:} Parents/students never informed of their rights (access, rectification, erasure).
                        \end{enumerate}
              \end{itemize}
        \item \textbf{Corrective Measures (Prioritised):}
              \begin{enumerate}
                  \item \textbf{Immediate:} Remove all grade PDFs and photos from WhatsApp groups. Delete the ``Anyone with link'' Drive share. Restrict Drive to \textit{``Restricted - Only people added can open''} with individual school email accounts.
                  \item \textbf{Short-term (1 month):}
                        \begin{itemize}
                            \item Deploy a \textbf{School Management System (SMS)} (e.g., \textit{OpenSIS}, \textit{SchoolTool}, or local Cameroonian EdTech like \textit{EduBridge}, \textit{Klasroom}) with role-based access: Parents see \textbf{only} their child's grades, attendance, fees. Teachers enter grades. Admin manages users.
                            \item Obtain \textbf{explicit written consent forms} for: (a) using student photos on school website/social media, (b) sharing parent contact info with class reps (opt-in).
                            \item Appoint a \textbf{DPO} (can be a trained teacher/admin) and register with ANTIC.
                            \item Publish a \textbf{Privacy Notice} on school website and in admission pack.
                        \end{itemize}
                  \item \textbf{Ongoing:}
                        \begin{itemize}
                            \item Train staff on data protection (don't send grades via WhatsApp/email).
                            \item Conduct a \textbf{Data Protection Impact Assessment (DPIA)} for any new system (e.g., biometric attendance).
                            \item Establish \textbf{Data Breach Response Plan}: detect $\rightarrow$ contain $\rightarrow$ assess risk $\rightarrow$ notify ANTIC (72h) $\rightarrow$ notify affected parents $\rightarrow$ document.
                        \end{itemize}
              \end{enumerate}
        \item \textbf{WhatsApp Alternative for Announcements:} Use \textbf{WhatsApp Broadcast List} (one-way, recipients don't see each other) or \textbf{WhatsApp Channel} (new feature) for general announcements. Never share personal data there.
    \end{enumerate}
    \textbf{Examiner's focus:} You must cite \textbf{Law No. 2024/014} (new, shows up-to-date knowledge) and \textbf{Law No. 2010/012} (cybercrime). Distinguish \textit{personal data} vs \textit{sensitive data}. Know the \textbf{72-hour breach notification} rule. Practical remedy = \textbf{role-based access system}, not just ``be careful''.
\end{exoresolu}

\begin{methode}[Evaluating the Credibility of Online Information (Fact-Checking)]
\begin{enumerate}
    \item \textbf{Source Evaluation (CRAAP Test adapted)}:
          \begin{itemize}
              \item \textbf{C}urrency: When published/updated? Is it current for the topic?
              \item \textbf{R}elevance: Does it answer your question? Audience level?
              \item \textbf{A}uthority: Who is the author/publisher? Credentials? Institutional affiliation (university, ministry, reputable NGO, established media)? In Cameroon: \textit{CRTV, Cameroon Tribune, ANTIC, MINSANTE, INS, WHO Africa} $>$ random blog/WhatsApp forward.
              \item \textbf{A}ccuracy: Evidence? Citations? Peer-reviewed? Corroborated by other independent sources?
              \item \textbf{P}urpose: Inform, persuade, sell, entertain, mislead? Bias? Sponsored content?
          \end{itemize}
    \item \textbf{Technical Verification Steps}:
          \begin{enumerate}
              \item \textbf{Reverse Image Search:} Google Lens, TinEye, Yandex Images $\rightarrow$ find original context/date of a photo/video.
              \item \textbf{Video Verification:} InVID/WeVerify plugin (browser) $\rightarrow$ keyframe extraction, metadata, geolocation (shadows, landmarks).
              \item \textbf{Domain Check:} Whois lookup (age, registrar), check for typosquatting (e.g., \texttt{antic-cm.org} vs \texttt{antic.cm}).
              \item \textbf{Fact-Checking Sites:} \textit{Africa Check}, \textit{Dubawa}, \textit{PesaCheck}, \textit{DataCheck Cameroon}, \textit{AFP Fact Check}.
              \item \textbf{Official Sources:} Go directly to ministry websites (\texttt{www.minsante.cm}, \texttt{www.minesec.cm}, \texttt{www.antic.cm}), not screenshots.
          \end{enumerate}
    \item \textbf{Red Flags for Misinformation in Cameroon Context}:
          \begin{itemize}
              \item \textit{``Forwarded many times''} on WhatsApp.
              \item Audio note from ``Dr. X at Yaoundé Central Hospital'' — no name, no verifiable ID.
              \item Screenshot of a tweet/post without link (easily faked).
              \item Urgent tone: \textit{``Share before they delete it!''}
              \item Ethnic/tribal/political inflammatory content near elections.
              \item Miracle cures (COVID, Ebola, cholera) — always check MINSANTE/WHO.
          \end{itemize}
    \item \textbf{Action:} \textbf{Stop} $\rightarrow$ \textbf{Think} $\rightarrow$ \textbf{Check} $\rightarrow$ \textbf{Share (only if verified)}. If unsure, \textbf{do not forward}.
\end{enumerate}
\end{methode}

\begin{exoresolu}[Fact-Checking a Viral WhatsApp Audio on Cholera Outbreak]
    \textbf{Statement:} You receive a WhatsApp voice note (2 min 13 sec) in a family group. A man claims to be ``Doctor Ngwa from Buea Regional Hospital''. He says: \textit{``Cholera is spreading fast in Buea and Limbe. The water from CAMWATER is contaminated. Do not drink tap water. Boil everything. The government is hiding the numbers. Share this to save lives.''} The message has the \textit{``Forwarded many times''} label. Apply a rigorous fact-checking methodology to determine the credibility and decide whether to forward.
    \tcblower
    \textbf{Detailed Solution:}
    \begin{enumerate}
        \item \textbf{STOP:} Do not forward. Acknowledge the fear appeal (urgency, authority impersonation, conspiracy).
        \item \textbf{Source Analysis (Authority):}
              \begin{itemize}
                  \item \textit{Claimed:} ``Dr. Ngwa, Buea Regional Hospital''.
                  \item \textit{Verification:} Search \texttt{``Dr Ngwa Buea Regional Hospital cholera''} on Google. Check hospital staff directory (if public). Call the hospital switchboard (\texttt{+237 233 32 2X XX}) — ask for the Director or PRO. \textbf{Result:} No Dr. Ngwa in infectious disease unit. The voice doesn't match known doctors.
              \end{itemize}
        \item \textbf{Cross-Check with Official Sources (Accuracy):}
              \begin{itemize}
                  \item \textbf{MINSANTE Website (\texttt{minsante.cm}) / Twitter (\texttt{@MinSanteCm}):} Check latest \textit{Communiqué} / \textit{Situation Report} on cholera. As of today, the official report shows \textbf{0 cases in Southwest Region} (hypothetical for example).
                  \item \textbf{WHO Africa / Cameroon:} \texttt{www.afro.who.int} $\rightarrow$ Disease Outbreak News.
                  \item \textbf{CAMWATER (\texttt{camwater.cm}):} Press release on water quality? None.
                  \item \textbf{ANTIC / ANTIC Fake News Alert:} Check \texttt{antic.cm} or their Facebook page for debunks.
              \end{itemize}
        \item \textbf{Content Analysis (Purpose / Red Flags):}
              \begin{itemize}
                  \item \textit{``Government is hiding numbers''} $\rightarrow$ Conspiracy narrative, undermines trust in institutions.
                  \item \textit{``Share to save lives''} $\rightarrow$ Emotional manipulation, virality mechanic.
                  \item \textit{No specific data:} No case numbers, dates, lab confirmation, ward names.
                  \item \textit{Audio only} $\rightarrow$ Cannot verify speaker visually; easy to record anonymously.
              \end{itemize}
        \item \textbf{Technical Check:}
              \begin{itemize}
                  \item Forwarded many times label $\rightarrow$ High probability of misinformation chain.
                  \item Search key phrases in quotes on Google: \textit{``CAMWATER contaminated cholera Buea''} $\rightarrow$ No credible news articles (Cameroon Tribune, CRTV, Journal du Cameroun, Actu Cameroun).
              \end{itemize}
        \item \textbf{Conclusion:} \textbf{UNVERIFIED / LIKELY FALSE.} The claim contradicts official health sources. The source is anonymous/unverifiable. The narrative matches known misinformation patterns.
        \item \textbf{Action:}
              \begin{enumerate}
                  \item \textbf{Do not forward.}
                  \item \textbf{Reply in group (calmly):} \textit{``Family, I checked with MINSANTE and Buea Regional Hospital. There's no cholera outbreak reported. The audio is unverified. Please don't spread unconfirmed health alerts — they cause panic. Official info: @MinSanteCm or minsante.cm.''}
                  \item \textbf{Report on WhatsApp:} Long-press $\rightarrow$ \textit{Report}.
                  \item \textbf{Share official link:} Post the latest MINSANTE cholera situation report link.
              \end{enumerate}
    \end{enumerate}
    \textbf{Key lesson:} In health emergencies, \textbf{only} trust official ministry channels and WHO. WhatsApp audio notes from ``doctors'' are the \#1 vector of health misinformation in Cameroon.
\end{exoresolu}

\begin{methode}[Leveraging Social Networks for Professional Development and Employability]
\begin{enumerate}
    \item \textbf{Build a Professional Personal Brand}:
          \begin{enumerate}
              \item \textbf{LinkedIn Profile Optimisation:} Professional photo (neutral background), Headline: \textit{``Final-year [Field] Student | [Key Skill 1] | [Key Skill 2] | Seeking Internship/Entry Role''}. About: 3-paragraph story (Who I am, What I do, What I want). Experience: Include \textbf{all} internships, volunteer work, leadership roles (class prefect, club president), projects — use action verbs + metrics (\textit{``Led team of 5 to organise career fair for 200+ students''}). Skills: Add 15+ relevant skills (endorsements). Recommendations: Request 2–3 from supervisors/lecturers.
              \item \textbf{Clean Digital Footprint:} (See Method 2) — Google yourself, remove/secure.
              \item \textbf{Content Strategy:} Post \textbf{once a week}: Share an article on your field with your insight (``I found this interesting because...''), comment thoughtfully on posts by target companies (Orange, MTN, banks, NGOs), share your project work (PDF on SlideShare/LinkedIn Articles).
          \end{enumerate}
    \item \textbf{Network Strategically}:
          \begin{itemize}
              \item \textbf{Alumni Search:} LinkedIn $\rightarrow$ School $\rightarrow$ Alumni $\rightarrow$ Filter by company/role $\rightarrow$ Connect with personalised note: \textit{``Hi [Name], I'm a final-year [Course] student at [School]. I admire your career at [Company]. May I ask one quick question about...?''}
              \item \textbf{Events:} Attend (physical/virtual) career fairs (e.g., \textit{Salon de l'Emploi}, \textit{YouthConnekt}), webinars. Connect with speakers within 24h.
              \item \textbf{WhatsApp/Telegram Professional Groups:} Join \textit{verified} groups for your field (e.g., \textit{Cameroon Software Engineers}, \textit{Young Accountants Cameroon}) — contribute, don't just lurk.
          \end{itemize}
    \item \textbf{Job Application via Social Media}:
          \begin{itemize}
              \item Follow target companies' \textbf{Careers Pages} (LinkedIn Jobs, company website, Facebook Page).
              \item Set \textbf{Job Alerts} on LinkedIn (keywords: \textit{Internship, Graduate Trainee, Junior} + location: \textit{Cameroon}).
              \item Apply \textbf{directly on company website} (ATS-friendly) — LinkedIn Easy Apply is convenient but less tailored.
              \item \textbf{Beware of Scams:} No legitimate employer asks for money (``registration fee'', ``medical fee'') before interview. Verify company registration (RCCM number) on \texttt{www.mcp.gov.cm} or \texttt{www.antic.cm}.
          \end{itemize}
    \item \textbf{Showcase Portfolio}:
          \begin{itemize}
              \item \textbf{Tech/Design:} GitHub (code), Behance/Dribbble (design), YouTube (demo videos).
              \item \textbf{Writing/Research:} LinkedIn Articles, Medium, ResearchGate, Google Scholar profile.
              \item \textbf{General:} Personal website (WordPress/Wix/Google Sites — free) with CV, projects, contact. Link in LinkedIn bio.
          \end{itemize}
\end{enumerate}
\end{methode}

\begin{exoresolu}[Building a LinkedIn Profile for a GCE Graduate Seeking an Apprenticeship]
    \textbf{Statement:} Ngwa, 19, just passed his GCE A-Level (Math, Physics, Chemistry) at GBHS Kumbo. He wants a Higher Technician apprenticeship (\textit{Brevet de Technicien Supérieur}) in \textit{Industrial Maintenance} at a company like CDC, SONARA, or a brewery (SABC/UCB). He has done a 2-month holiday internship at a local garage (welding, lathe) and was class prefect. Write the exact \textbf{Headline}, \textbf{About} section (3 paragraphs), and \textbf{Experience} entries he should put on LinkedIn. List 10 strategic skills to add.
    \tcblower
    \textbf{Detailed Solution:}
    \textbf{Headline:}
    \texttt{Final-Year GCE A-Level Graduate (Math/Physics/Chemistry) | Aspiring Industrial Maintenance Technician | Lathe \& Welding Experience | Seeking BTS Apprenticeship (CDC, SONARA, SABC)}

    \textbf{About:}
    \begin{quote}
    \textbf{Paragraph 1 (Who I am):} I am a disciplined and hands-on recent GCE Advanced Level graduate (Math, Physics, Chemistry) from Government Bilingual High School Kumbo, with a passion for keeping industrial machinery running efficiently. My goal is to secure a Brevet de Technicien Supérieur (BTS) apprenticeship in Industrial Maintenance starting September 2025, where I can apply my strong theoretical foundation and practical workshop experience.

    \textbf{Paragraph 2 (What I've done):} As Class Prefect (2023–2024), I led a team of 12 prefects managing discipline and events for 300+ students, honing my leadership, communication, and problem-solving under pressure. During the 2024 long holidays, I completed a 2-month voluntary internship at \textit{Ngwa's Auto \& Metal Workshop} in Kumbo, where I gained practical skills in: manual lathe operation (turning, facing, threading), arc/MIG welding (joint preparation, safety), preventive maintenance of generators and pumps, and reading basic technical drawings. I also hold a valid \textbf{First Aid Certificate (Red Cross Cameroon)}.

    \textbf{Paragraph 3 (What I want + Call to Action):} I am eager to bring my work ethic, safety consciousness, and eagerness to learn to a structured industrial environment like CDC, SONARA, SABC, or a similar manufacturing plant. I am available for interviews and aptitude tests immediately. Open to relocating within Cameroon. \textbf{Let's connect — I welcome advice from maintenance engineers and HR professionals in the sector.}
    \end{quote}

    \textbf{Experience Entries:}
    \begin{enumerate}
        \item \textbf{Class Prefect} — Government Bilingual High School Kumbo \hfill \textit{Sep 2023 – Jun 2024}
              \begin{itemize}
                  \item Led a prefectural board of 12, coordinating daily discipline, assembly logistics, and inter-house sports for 300+ students.
                  \item Mediated 15+ student conflicts per term, reducing escalation to administration by 40\%.
                  \item Organised the 2024 End-of-Year Cultural Festival (budget 150,000 FCFA, 500 attendees) — managed vendor contracts, security, and cleanup.
              \end{itemize}
        \item \textbf{Workshop Intern (Voluntary)} — Ngwa's Auto \& Metal Workshop, Kumbo \hfill \textit{Jul – Aug 2024}
              \begin{itemize}
                  \item Operated manual centre lathe for turning, facing, and thread cutting (M6–M20) under supervision; produced 20+ replacement shafts/bushings for water pumps.
                  \item Performed arc welding (SMAW) on mild steel frames (gate repairs, trailer reinforcements); followed PPE and fire safety protocols.
                  \item Assisted in preventive maintenance of 5kVA–20kVA generators: oil changes, filter replacement, battery testing, load bank testing.
                  \item Interpreted basic mechanical drawings (ISO views, tolerances) to fabricate parts.
              \end{itemize}
    \end{enumerate}

    \textbf{Top 10 Skills to Add (prioritised for Industrial Maintenance):}
    \begin{enumerate}
        \item Preventive Maintenance
        \item Manual Lathe Operation
        \item Arc Welding (SMAW)
        \item Technical Drawing Interpretation
        \item Industrial Safety (HSE)
        \item Troubleshooting (Mechanical)
        \item Hand \& Power Tools
        \item Teamwork \& Communication
        \item Time Management
        \item First Aid / CPR
    \end{enumerate}

    \textbf{Next Steps for Ngwa:}
    \begin{itemize}
        \item Upload a PDF of his \textit{GCE Results Slip} and \textit{Internship Attestation} to the \textit{Featured} section.
        \item Connect with \textbf{HR Managers} at CDC (\texttt{cdc-cameroon.com}), SONARA, SABC — search ``Human Resources'', ``Maintenance Manager''.
        \item Join LinkedIn Groups: \textit{``Cameroon Engineers' Forum''}, \textit{``Maintenance Professionals Africa''}.
        \item Set Job Alert: \textit{``Industrial Maintenance Apprentice''}, \textit{``BTS Maintenance''}, \textit{``Technicien Maintenance''} — Location: \textit{Cameroon}.
    \end{itemize}
    \textbf{Examiner's note:} A GCE-level student \textbf{can} have a strong LinkedIn profile by framing leadership (prefect), voluntary practical experience, and clear career intent. The headline must contain \textbf{keywords} recruiters search for.
\end{exoresolu}

\newpage

\coursec{Exercises}

\courssub{I check my knowledge}

\begin{exercice}[MCQ: Social network classification]
For each platform below, choose the \textbf{single best category} from the list: \emph{General social network}, \emph{Professional network}, \emph{Instant messaging}, \emph{Collaborative encyclopedia}, \emph{Video conferencing}, \emph{Short-form video sharing}.
\begin{enumerate}
    \item Facebook
    \item LinkedIn
    \item WhatsApp
    \item Wikipedia
    \item Zoom
    \item TikTok
\end{enumerate}
\end{exercice}

\begin{exercice}[True/False with justification]
State whether each statement is \textbf{True} or \textbf{False}. Justify your answer in one sentence.
\begin{enumerate}
    \item A post shared with the audience setting ``Friends'' on Facebook can never be seen by people outside that friend list.
    \item Forwarding a message containing false information in a WhatsApp group has no legal consequence in Cameroon.
    \item Two-factor authentication (2FA) prevents anyone from accessing your account even if they know your password.
    \item A digital footprint only includes information you deliberately publish online.
\end{enumerate}
\end{exercice}

\begin{exercice}[Definitions]
Define the following terms clearly and concisely (one to two sentences each).
\begin{enumerate}
    \item Hashtag
    \item News feed
    \item Wiki
    \item Netiquette
    \item Digital footprint
\end{enumerate}
\end{exercice}

\begin{exercice}[Direct calculation: Password entropy]
A student creates a password using only lowercase letters (26 possibilities per character).
\begin{enumerate}
    \item Calculate the number of possible passwords of length 8.
    \item If the student adds one uppercase letter, one digit, and one special character (total 26+26+10+10 = 72 possibilities per character), how many times larger is the search space for an 8-character password compared to the lowercase-only version? Express your answer as a power of 10 (approximate).
\end{enumerate}
\end{exercice}

\courssub{I practise}

\begin{exercice}[Account security scenario]
A Form 5 student discovers that their Facebook account has been hacked. The attacker changed the password, recovery email, and phone number, then posted a message asking friends to send money via Mobile Money (MTN MoMo / Orange Money) to a specific number.
\begin{enumerate}
    \item List \textbf{five preventive steps} the student should have taken \emph{before} the hack to secure the account.
    \item List \textbf{four immediate actions} the student must take \emph{after} discovering the hack.
\end{enumerate}
\end{exercice}

\begin{exercice}[Collaborative editing vs. email]
A group of five Upper Sixth students must write a 15-page research report together. They consider two workflows:
\begin{itemize}
    \item \textbf{Workflow A:} One student creates a Google Doc, shares it with the other four (edit access), and they all work simultaneously.
    \item \textbf{Workflow B:} The student emails a Word file to the others; each edits their own copy and emails it back; one student merges all changes manually.
\end{itemize}
Explain why Workflow A is superior for this task. Cite \textbf{at least three specific features} of collaborative editing platforms (e.g., real-time co-editing, version history, commenting, suggestion mode, access control) and describe how each feature solves a concrete problem that occurs in Workflow B.
\end{exercice}

\begin{exercice}[Information verification case study]
A photo circulates in a WhatsApp group claiming: ``\emph{The GCE Board has postponed the Advanced Level examinations by two weeks. New timetable attached.}'' The photo shows a document with the GCE Board logo and a signature.
Describe, \textbf{step by step}, how a responsible student should verify this information before sharing it further. Your answer must include at least \textbf{five distinct verification actions}, referencing official Cameroonian sources where appropriate.
\end{exercice}

\begin{exercice}[Matching exercise]
Match each term in Column A with its correct definition in Column B. Write the pairs (e.g., 1--f, 2--a, ...).
\begin{center}
\begin{tabular}{|c|l|}\hline
\textbf{Column A} & \textbf{Column B} \\ \hline
1. Hashtag & a. A set of rules for polite and responsible behaviour online \\ \hline
2. News feed & b. A word or phrase preceded by \# used to categorise content \\ \hline
3. Wiki & c. The trail of data you leave when using the internet \\ \hline
4. 2FA (Two-Factor Authentication) & d. A website that allows collaborative editing of its content \\ \hline
5. Netiquette & e. A security method requiring two different authentication factors \\ \hline
6. Digital footprint & f. The constantly updating list of posts seen on a social network home page \\ \hline
\end{tabular}
\end{center}
\end{exercice}

\courssub{I prepare for the GCE}

\begin{exercice}[{Essay: Impact of social networks on Cameroonian students] [15 marks}]
\textbf{``Social networks do more good than harm to Cameroonian students.''}

Discuss this statement. Your essay must:
\begin{itemize}
    \item Present \textbf{at least three arguments supporting} the view that social networks are beneficial (e.g., access to educational resources, networking, civic engagement, digital skills).
    \item Present \textbf{at least three arguments supporting} the view that social networks are harmful (e.g., distraction, cyberbullying, misinformation, privacy risks, mental health).
    \item Provide a \textbf{personal, justified conclusion} that weighs the arguments and states your position clearly.
\end{itemize}
Structure your answer with clear paragraphs. Use relevant Cameroonian examples (e.g., WhatsApp study groups, Facebook scholarship alerts, TikTok trends, ANTIC campaigns).
\end{exercice}

\begin{exercice}[{Structured question: Privacy settings checklist] [10 marks}]
You are asked to design a \textbf{Privacy Settings Checklist} for a Form 5 student opening their first social network account (e.g., Facebook, Instagram, TikTok).
\begin{enumerate}
    \item List \textbf{eight essential privacy/settings actions} the student should configure immediately after creating the account. For each action, give a one-sentence explanation of \emph{why} it matters. [8 marks]
    \item State \textbf{two risks} the student faces if they leave the account on default public settings. [2 marks]
\end{enumerate}
Present your answer as a numbered checklist.
\end{exercice}

\begin{exercice}[{Case study: Collaborative project management] [10 marks}]
A team of four Upper Sixth students is building a website for their school's 50th anniversary. They decide to use GitHub for code version control, Google Drive for shared assets (images, documents), and a WhatsApp group for quick communication.
\begin{enumerate}
    \item Explain \textbf{why GitHub is more suitable than Google Drive} for managing the website's source code (HTML, CSS, JavaScript). Mention at least two specific Git/GitHub features. [4 marks]
    \item The team leader wants to enforce a rule: ``No direct commits to the \texttt{main} branch; all changes must go through a Pull Request.'' Explain \textbf{two advantages} of this workflow for a student team. [3 marks]
    \item Identify \textbf{one limitation} of using WhatsApp for project coordination and propose a better alternative tool (name it and justify). [3 marks]
\end{enumerate}
\end{exercice}

\newpage

\coursec{Solutions to the exercises}

\begin{corrige}[Exercise 1 — MCQ: Social network classification]
\begin{enumerate}
    \item \textbf{Facebook — General social network.} It connects people for personal and social purposes (friends, family, groups, pages) and is not specialised in one professional or media function.
    \item \textbf{LinkedIn — Professional network.} Its core purpose is career development: CVs, job offers, business contacts and professional content.
    \item \textbf{WhatsApp — Instant messaging.} Although it has groups and status updates, its primary function is real-time one-to-one and group messaging, voice and video calls.
    \item \textbf{Wikipedia — Collaborative encyclopedia.} Its content is written and edited collectively by volunteers using wiki technology; it is a \emph{collaborative environment}, not a social network in the usual sense.
    \item \textbf{Zoom — Video conferencing.} It is designed for live audio-video meetings, webinars and online classes.
    \item \textbf{TikTok — Short-form video sharing.} Its defining feature is the creation and consumption of short videos, driven by a recommendation algorithm.
\end{enumerate}
\begin{attention}
A frequent mistake is to classify WhatsApp as a ``general social network''. In the GCE classification, a platform is categorised by its \emph{dominant function}: WhatsApp's dominant function is messaging, even though it includes some social features like status updates and groups.
\end{attention}
\end{corrige}

\begin{corrige}[Exercise 2 — True/False with justification]
\begin{enumerate}
    \item \textbf{False.} A ``Friends'' post can still leak beyond the intended audience: a friend can take a screenshot, copy the text, or share the content manually, and friends of friends may see it if someone is tagged. Audience settings reduce visibility but do not guarantee absolute confidentiality.
    \item \textbf{False.} In Cameroon, the 2010 Law on Cybersecurity and Cybercriminality (Law No. 2010/012 of 21 December 2010) punishes the publication and propagation of false information through electronic communications; forwarding fake news can therefore have legal consequences, and ANTIC regularly sensitises the public on this.
    \item \textbf{False.} 2FA makes unauthorised access \emph{much harder} because a second factor (code by SMS or authenticator app) is required, but it does not make access impossible: SIM-swap attacks, phishing of codes, or theft of the second device can defeat it. No security measure is absolute.
    \item \textbf{False.} The digital footprint also includes \emph{passive} data collected without deliberate publication: browsing history, cookies, location data, likes, and metadata generated simply by using online services.
\end{enumerate}
\end{corrige}

\begin{corrige}[Exercise 3 — Definitions]
\begin{enumerate}
    \item \textbf{Hashtag:} a keyword or phrase preceded by the symbol \# (e.g., \#GCE2025) used to label and group publications by topic, making them easier to search and follow on social networks.
    \item \textbf{News feed:} the continuously updated stream of posts, photos, videos and announcements displayed on a user's home page, selected and ordered by the platform's algorithm according to the user's connections and activity.
    \item \textbf{Wiki:} a website whose pages can be created and modified directly and collaboratively by its users through a web browser, usually keeping a history of all modifications (example: Wikipedia).
    \item \textbf{Netiquette:} the set of rules of polite, respectful and responsible behaviour to observe when communicating on the internet (no insults, no spam, respect of privacy, writing clearly, not typing in capital letters, etc.).
    \item \textbf{Digital footprint:} the complete trail of data a person leaves behind when using the internet, including deliberately published content (active footprint) and data collected automatically such as cookies and location (passive footprint).
\end{enumerate}
\end{corrige}

\begin{corrige}[Exercise 4 — Direct calculation: Password entropy]
\textbf{1.} Each of the 8 positions can be filled independently by any of the 26 lowercase letters. By the multiplicative principle:
\[
N = 26^8.
\]
Computing step by step: $26^2 = 676$, $26^4 = 676^2 = 456\,976$, $26^8 = 456\,976^2 = 208\,827\,064\,576$.
\[
\boxed{N = 26^8 \approx 2.09 \times 10^{11}\ \text{passwords}.}
\]

\textbf{2.} With 72 possibilities per character (26 lowercase + 26 uppercase + 10 digits + 10 common symbols), the new search space is $72^8$. The ratio is:
\[
\frac{72^8}{26^8} = \left(\frac{72}{26}\right)^8.
\]
Now $\dfrac{72}{26} = \dfrac{36}{13} \approx 2.76923$, so
\[
\left(\frac{72}{26}\right)^8 \approx 2.76923^8.
\]
Using logarithms: $8 \log_{10}(2.76923) \approx 8 \times 0.4424 \approx 3.539$, hence
\[
\left(\frac{72}{26}\right)^8 \approx 10^{3.54} \approx 3.5 \times 10^{3}.
\]
\[
\boxed{\text{The search space is about } 10^{3.5} \approx 3\,500 \text{ times larger.}}
\]
\textbf{Interpretation:} diversifying the character set multiplies the attacker's work by several thousand for the same password length, which is why mixing cases, digits and symbols is strongly recommended.
\end{corrige}

\begin{corrige}[Exercise 5 — Account security scenario]
\textbf{1. Five preventive steps (before the hack):}
\begin{enumerate}
    \item \textbf{Use a strong, unique password} (at least 12 characters mixing uppercase, lowercase, digits and symbols), never reused on other sites, so that a leak elsewhere cannot compromise Facebook.
    \item \textbf{Enable two-factor authentication (2FA)}, preferably with an authenticator app (Google Authenticator, Microsoft Authenticator), so that knowing the password alone is insufficient to log in.
    \item \textbf{Activate login alerts} to be notified immediately of any connection from an unknown device or location.
    \item \textbf{Beware of phishing}: never enter credentials through links received by SMS, email or Messenger; always check the address \texttt{facebook.com} before typing a password.
    \item \textbf{Keep recovery information up to date and protected} (secure the recovery email with its own strong password and 2FA; never share codes received by SMS, since MoMo-style scams often start by tricking victims into revealing verification codes).
\end{enumerate}

\textbf{2. Four immediate actions (after discovering the hack):}
\begin{enumerate}
    \item \textbf{Start Facebook's account recovery procedure} (\texttt{facebook.com/hacked}) to prove identity and regain control of the account.
    \item \textbf{Warn all contacts immediately} (by phone, SMS, or another account) that the money request is a scam, so that nobody sends MTN MoMo or Orange Money to the attacker's number.
    \item \textbf{Report the fraudulent Mobile Money number} to MTN or Orange customer service so the wallet can be flagged or blocked, and report the compromised account to Facebook.
    \item \textbf{Change the passwords of all linked accounts} (especially the recovery email) and, once access is restored, review active sessions, remove unknown devices, and enable 2FA. A complaint can also be lodged with the police or reported to ANTIC.
\end{enumerate}
\end{corrige}

\begin{corrige}[Exercise 6 — Collaborative editing vs. email]
Workflow A is superior because a collaborative editing platform is designed precisely to eliminate the problems created by circulating file copies by email.

\textbf{1. Real-time co-editing.} In Workflow B, only one person can work on a given copy at a time, and the group must wait for each round of emails; with five students, several conflicting versions quickly exist and nobody knows which is current. In Workflow A, all five students edit the \emph{same single document} simultaneously and see each other's changes live, so there is always exactly one up-to-date version.

\textbf{2. Version history.} In Workflow B, if the student merging the files accidentally deletes a section or merges wrongly, the lost content may be unrecoverable. Google Docs keeps an automatic history of every version, so the group can see who changed what and restore any previous state in one click.

\textbf{3. Commenting and suggestion mode.} In Workflow B, remarks are written inside the text or in separate emails, and the merger must interpret and apply them by hand — a source of errors and omissions. In Workflow A, a member can attach a comment to an exact passage or propose a change in suggestion mode, which the group accepts or rejects without altering the original text until decided.

\textbf{4. Access control.} In Workflow B, anyone who forwards the email can modify or leak the file. In Workflow A, the owner grants precise rights (view / comment / edit) to named accounts and can revoke them at any time.

\textbf{Conclusion:} Workflow A saves time, prevents conflicting copies, protects the work through history and access rights, and structures feedback — all of which Workflow B cannot guarantee.
\end{corrige}

\begin{corrige}[Exercise 7 — Information verification case study]
A responsible student should apply the \emph{verify before you share} reflex as follows:
\begin{enumerate}
    \item \textbf{Do not forward immediately.} The first action is to stop the chain: an unverified forward makes the student a spreader of possible fake news, which is punishable under Cameroon's 2010 law on cybersecurity and cybercriminality (Law No. 2010/012).
    \item \textbf{Check the official GCE Board sources.} Visit the official website of the Cameroon GCE Board (\texttt{www.cameroongceboard.org}) and its official communication channels to see whether any communiqué about a postponement has been published. An authentic decision of this importance always appears there.
    \item \textbf{Check the Ministry of Secondary Education (MINESEC)} official website (\texttt{www.minesec.gov.cm}) and verified social media pages, since the GCE Board operates under its authority and a calendar change would be confirmed there.
    \item \textbf{Examine the document critically.} Look for signs of forgery: poor-quality logo, spelling or grammar errors, wrong date format, inconsistent signature, absence of an official reference number — fake documents circulating on WhatsApp often contain such defects.
    \item \textbf{Cross-check with credible media.} Consult trusted national media (e.g., CRTV, Cameroon Tribune) or established news sites: a two-week postponement of a national examination would be widely reported.
    \item \textbf{Ask an authoritative person.} Consult a teacher, the school administration, or the principal, who receive official correspondence from the GCE Board through official channels.
    \item \textbf{Use fact-checking reflexes.} Search the exact claim on a search engine to see whether fact-checkers or the GCE Board itself have already denied it (institutions sometimes publish ``fake news'' alerts).
\end{enumerate}
\textbf{Conclusion:} only after confirmation from at least one \emph{official} source (GCE Board or MINESEC) should the information be considered reliable — and even then, sharing the official communiqué itself is preferable to forwarding the suspicious photo.
\end{corrige}

\begin{corrige}[Exercise 8 — Matching exercise]
The correct pairs are:
\begin{center}
\begin{tabular}{|c|l|}\hline
\rowcolor{popblueL} \textbf{Pair} & \textbf{Justification} \\ \hline
1 -- b & A hashtag is precisely a word or phrase preceded by \# used to categorise content. \\ \hline
2 -- f & The news feed is the constantly updating list of posts on the home page. \\ \hline
3 -- d & A wiki is a website allowing collaborative editing of its content. \\ \hline
4 -- e & 2FA requires two different authentication factors (e.g., password + SMS code). \\ \hline
5 -- a & Netiquette is the set of rules for polite and responsible behaviour online. \\ \hline
6 -- c & The digital footprint is the trail of data left when using the internet. \\ \hline
\end{tabular}
\end{center}
\textbf{Final answer:} 1--b, 2--f, 3--d, 4--e, 5--a, 6--c.
\end{corrige}

\begin{corrige}[Exercise 9 — Essay: Impact of social networks on Cameroonian students (15 marks)]
\textbf{Examiner expectations:} a structured essay (introduction, balanced development, personal conclusion), at least three arguments on each side, Cameroonian illustrations, correct language and a clearly justified personal position.

\textbf{Indicative mark scheme (15 marks):}
\begin{center}
\begin{tabularx}{\linewidth}{|l|X|c|}
\hline
\rowcolor{popblueL} \textbf{Component} & \textbf{Expectation} & \textbf{Marks} \\ \hline
Introduction & Statement of the issue and announcement of the plan & 2 \\ \hline
Benefits & At least 3 developed arguments with examples & 4 \\ \hline
Harms & At least 3 developed arguments with examples & 4 \\ \hline
Conclusion & Personal, justified position weighing both sides & 3 \\ \hline
Quality & Structure, coherence, spelling, Cameroonian relevance & 2 \\ \hline
\end{tabularx}
\end{center}

\textbf{Model outline and content:}

\emph{Introduction.} Social networks such as WhatsApp, Facebook and TikTok have become part of the daily life of Cameroonian students, from Yaoundé to Bamenda. Whether they do more good than harm is a debated question that deserves a balanced examination.

\emph{Arguments for the benefits (at least three):}
\begin{itemize}
    \item \textbf{Access to educational resources:} WhatsApp study groups allow classmates to share past GCE papers, lessons and explanations; YouTube and Facebook pages offer free tutorials in mathematics, physics and ICT.
    \item \textbf{Information and opportunities:} Facebook pages and groups publish scholarship alerts, competition announcements (concours) and GCE Board information faster than traditional channels.
    \item \textbf{Networking and digital skills:} by creating content and collaborating online, students acquire digital skills valued in the modern job market; platforms like LinkedIn connect advanced students with mentors and employers.
    \item \textbf{Civic engagement:} youths follow ANTIC sensitisation campaigns and public debates, developing civic awareness.
\end{itemize}

\emph{Arguments for the harms (at least three):}
\begin{itemize}
    \item \textbf{Distraction and poor results:} endless scrolling on TikTok or late-night chatting reduces study time and concentration, contributing to declining performance at school.
    \item \textbf{Misinformation:} fake news (e.g., false GCE postponements, miracle cures) spreads rapidly in WhatsApp groups and can mislead students who forward without verifying.
    \item \textbf{Cyberbullying and mental health:} mockery, humiliating photos and comparison with idealised lives cause stress, anxiety and loss of self-esteem.
    \item \textbf{Privacy and security risks:} oversharing personal data exposes students to hacking, identity theft and MoMo scams, as seen in account-hijacking cases.
\end{itemize}

\emph{Personal conclusion (example).} Social networks are a \emph{tool}: their effect depends on how they are used. For a disciplined student who verifies information, protects privacy and limits screen time, the educational and social benefits outweigh the harms. Therefore, rather than banning them, schools and families should teach digital literacy and netiquette — as ANTIC campaigns encourage — so that Cameroonian students can exploit the good while avoiding the traps.

\textbf{Examiner's advice:} an essay that presents only one side, or lists arguments without examples and without a personal justified conclusion, cannot earn full marks.
\end{corrige}

\begin{corrige}[Exercise 10 — Structured question: Privacy settings checklist (10 marks)]
\textbf{1. Eight essential privacy actions [8 marks — 1 mark per action with valid justification]:}
\begin{enumerate}
    \item \textbf{Set the profile to private / restrict the default audience} — so that only approved friends, and not the whole internet, can see personal posts and photos.
    \item \textbf{Create a strong, unique password} — to prevent easy guessing or reuse of a password leaked from another site.
    \item \textbf{Enable two-factor authentication (2FA)} — so that even a stolen password is not enough to enter the account.
    \item \textbf{Limit who can send friend/follow requests and messages} — to reduce contact from strangers, scammers and cyberbullies.
    \item \textbf{Control tagging and mentions (require approval before a tag appears)} — to stop others from associating the student's name with embarrassing or unwanted content.
    \item \textbf{Hide personal information (phone number, email, date of birth, school, location)} — because such data can be exploited for identity theft, stalking or targeted scams.
    \item \textbf{Disable or restrict location sharing} — so that the student's real-time whereabouts are not exposed to strangers.
    \item \textbf{Review connected apps and active sessions, and activate login alerts} — to detect and remove unknown devices or third-party apps abusing the account.
\end{enumerate}

\textbf{2. Two risks of default public settings [2 marks — 1 mark each]:}
\begin{itemize}
    \item \textbf{Exposure to strangers and predators:} anyone can view photos, personal details and location, facilitating cyberbullying, stalking or identity theft.
    \item \textbf{Permanent damage to reputation and digital footprint:} public posts can be copied, screenshotted and redistributed, harming future school, scholarship or job opportunities.
\end{itemize}
\end{corrige}

\begin{corrige}[Exercise 11 — Case study: Collaborative project management (10 marks)]
\textbf{1. Why GitHub is more suitable than Google Drive for source code [4 marks]:}

Google Drive stores files but does not understand code evolution. GitHub, built on Git, offers:
\begin{itemize}
    \item \textbf{Version control with commits and full history:} every change to the HTML/CSS/JavaScript files is recorded with an author, date and message; the team can compare versions and revert to any previous state if a change breaks the site — impossible with simple file copies on Drive.
    \item \textbf{Branching and merging:} each student can develop a feature on a separate branch without disturbing the working version, then merge it; Git detects and helps resolve conflicts when two members edit the same lines — whereas Drive would simply keep the last uploaded copy and silently overwrite the other.
\end{itemize}
(Also acceptable: issue tracking to assign bugs/tasks, and the distributed nature of Git giving every member a complete local backup of the project.)

\emph{Marking: 2 marks per feature clearly named and linked to a concrete advantage over Drive.}

\textbf{2. Two advantages of the Pull Request workflow [3 marks]:}
\begin{itemize}
    \item \textbf{Code review and quality:} before merging, teammates read and comment on the proposed changes, so errors, bugs and bad practices are caught early and everyone learns from the feedback.
    \item \textbf{Protection and traceability of the main branch:} \texttt{main} always contains a stable, working version of the site, and the history shows exactly who proposed, reviewed and approved each change — essential for a team where members have different skill levels.
\end{itemize}
(Also acceptable: it prevents accidental or untested direct commits and creates a natural discussion space per change.)

\emph{Marking: 1.5 marks per advantage well explained.}

\textbf{3. Limitation of WhatsApp and better alternative [3 marks]:}

\textbf{Limitation (1 mark):} WhatsApp mixes project discussion with personal chat, has no task tracking, and important decisions or files get lost in the message flow; it is not searchable or structured by topic.

\textbf{Alternative (2 marks):} a tool such as \textbf{Trello} (or GitHub Projects / Slack) is better: Trello organises work into boards, lists and cards, so each task (e.g., ``design homepage'', ``fix CSS bug'') can be assigned to a member, given a deadline and moved from ``To do'' to ``Done'', giving the whole team a clear, permanent view of progress — something a chat stream cannot provide.
\end{corrige}

\newpage

\coursec{Summary sheet}

\begin{popfigure}
\begin{tikzpicture}[font=\small,
central/.style={draw=popink, fill=popink, text=white, rounded corners=3pt, align=center, font=\bfseries, text width=3.6cm},
branch/.style={draw, rounded corners=3pt, align=center, text width=2.9cm, font=\scriptsize},
link/.style={draw=popmuted, thick}]
\node[central] (c) at (0,0) {Social Networks \& Collaborative Environments};
\node[branch, draw=popblue, fill=popblueL, align=center] (a) at (-4.6,2.6){\textbf{Concepts \& Types}\\ general, professional, media-sharing, microblogging};
\node[branch, draw=poporange, fill=popblueL, align=center] (b) at (0,3.4){\textbf{Key Actors}\\ users, providers, advertisers, regulators (ANTIC)};
\node[branch, draw=popgreen, fill=popgreenL, align=center] (d) at (4.6,2.6){\textbf{Using a Network}\\ account, profile, privacy settings, 2FA};
\node[branch, draw=poppurple, fill=popblueL, align=center] (e) at (4.6,-2.6){\textbf{Collaborative Tools}\\ cloud, wikis, forums, shared documents};
\node[branch, draw=popgold, fill=popgreenL, align=center] (f) at (0,-3.4){\textbf{Risks}\\ cyberbullying, addiction, data theft, fake news};
\node[branch, draw=poppink, fill=popblueL, align=center] (g) at (-4.6,-2.6){\textbf{Ethics \& Law}\\ netiquette, Law No.\ 2010/012, copyright};
\draw[link] (c) -- (a); \draw[link] (c) -- (b); \draw[link] (c) -- (d);
\draw[link] (c) -- (e); \draw[link] (c) -- (f); \draw[link] (c) -- (g);
\end{tikzpicture}
\legende{Overview of the chapter: social networks and collaborative environments.}
\end{popfigure}

\begin{aretenir}[Key definitions]
\begin{itemize}
\item \cle{Social network}: an online platform that allows users to create a public or semi-public profile, connect with other users and share content (text, photos, videos). Examples: Facebook, WhatsApp, X (formerly Twitter), TikTok, LinkedIn.
\item \cle{Collaborative environment}: a digital workspace where several users create, edit and share documents or projects together, often in real time. Examples: Google Workspace, Microsoft Teams, Moodle, wikis.
\item \cle{Netiquette}: the set of polite and responsible rules of behaviour on the Internet (respect, no spam, no insults, citing sources).
\item \cle{Digital identity}: all the information about a person available online (profiles, posts, photos, comments).
\item \cle{Privacy settings}: options controlling who can see your profile, posts and personal data.
\end{itemize}
\end{aretenir}

\begin{aretenir}[Types of social networks]
\begin{itemize}
\item \cle{General / friendship networks}: Facebook, WhatsApp --- connect friends and family.
\item \cle{Professional networks}: LinkedIn --- jobs, CVs, business contacts.
\item \cle{Media-sharing networks}: YouTube, TikTok, Instagram --- photos and videos.
\item \cle{Microblogging}: X (formerly Twitter) --- short public messages.
\end{itemize}
\textbf{Key actors}: users, platform providers (Meta, Google), advertisers, and regulators. In Cameroon, \cle{ANTIC} (National Agency for Information and Communication Technologies) promotes and monitors the secure use of ICT; operators such as MTN and Orange provide Internet access.
\end{aretenir}

\begin{methode}[Using a social network safely]
\begin{enumerate}
\item Create an account with a valid e-mail address and a \cle{strong password} (at least 8 characters, mixing upper- and lower-case letters, digits and symbols).
\item Complete the profile, then immediately configure the \cle{privacy settings} (who can see your posts, send you friend requests, or access your location).
\item Enable \cle{two-factor authentication} (2FA) where available, so that a code sent to your phone is required in addition to the password.
\item Publish responsibly: check information before sharing it, and never post other people's photos without their consent.
\end{enumerate}
\end{methode}

\begin{aretenir}[Collaborative tools to know]
\begin{itemize}
\item \cle{Cloud storage and sharing}: Google Drive, OneDrive, Dropbox.
\item \cle{Real-time co-editing}: Google Docs, Microsoft 365.
\item \cle{Communication}: e-mail, forums, videoconferencing (Zoom, Google Meet), instant messaging.
\item \cle{E-learning platforms}: Moodle, Google Classroom.
\end{itemize}
\textbf{Advantages}: teamwork at a distance, version history, time saving. \textbf{Limits}: need for a reliable Internet connection, possible conflicts during simultaneous editing, dependence on the service provider.
\end{aretenir}

\begin{aretenir}[Risks and protection]
\begin{itemize}
\item \textbf{Risks}: \cle{cyberbullying}, identity theft, scams and phishing, addiction, fake news, exposure of personal data.
\item \textbf{Protection}: strong and distinct passwords, limited personal information online, critical verification of sources, reporting and blocking abusers, balanced screen time.
\end{itemize}
\textbf{Legal framework (Cameroon)}: Law No.\ 2010/012 of 21 December 2010 on cybersecurity and cybercriminality in Cameroon punishes online offences such as identity theft, defamation and the publication of false information; ANTIC oversees digital trust. Copyright must be respected when sharing content.
\end{aretenir}

\begin{attention}[Common mistakes to avoid]
\begin{itemize}
\item Confusing a social network (relationships, profiles) with a collaborative tool (joint production of work).
\item Using the same weak password everywhere, or sharing passwords with friends.
\item Believing that deleting a post erases it completely: screenshots and archives may remain.
\item Publishing personal data (address, phone number, exam results) publicly.
\item Forwarding unverified information --- in Cameroon this can constitute an offence.
\end{itemize}
\end{attention}

\begin{aretenir}[GCE examination tips]
\begin{itemize}
\item Define terms precisely and always give one example (e.g.\ ``LinkedIn is a professional social network'').
\item For ``advantages/disadvantages'' questions, present balanced answers in two clear lists.
\item Quote Law No.\ 2010/012 on cybersecurity and cybercriminality, and ANTIC, when asked about regulation in Cameroon.
\item Practical questions may ask you to describe steps --- account creation, privacy configuration, sharing a document --- so answer in clear numbered steps.
\end{itemize}
\end{aretenir}

\findecours

\end{document}
